% =====================================================================
%  Anatomía de un LLM — Fase VI — Hito 6 (cierre del curso)
%  Examen de hito y proyecto final: diseño y construcción de optimizadores con IA
%  Compilar desde esta carpeta: latexmk -r ../latexmkrc H6-examen.tex
% =====================================================================
\documentclass[11pt,a4paper,twoside,openany]{book}
\usepackage{../estilo-apuntes}
\usepackage{makecell}
\usepackage{float}

\newcommand{\ApuntesCurso}{Anatomía de un LLM}
\newcommand{\ApuntesModulo}{Hito 6}

\hypersetup{pdftitle={Hito 6. Examen de la fase VI y proyecto final del curso},pdfauthor={Anatomía de un LLM}}
\makeatletter\g@addto@macro\UrlBreaks{\do\-\do\/}\makeatother
% Anchura del número de capítulo en el índice, para los apéndices III y IV
\makeatletter\patchcmd{\l@chapter}{1.5em}{2.4em}{}{}\makeatother

% --- Notación de la fase VI (la misma que en los apuntes M24–M28) ---------
\newcommand{\T}{^{\top}}
\DeclareMathOperator*{\argmin}{arg\,min}
\DeclareMathOperator*{\argmax}{arg\,max}
\DeclareMathOperator{\clip}{clip}
\newcommand{\Prob}{\mathbb{P}}
\newcommand{\N}{\mathcal{N}}
\newcommand{\EI}{\mathrm{EI}}
\newcommand{\PI}{\mathrm{PI}}
\newcommand{\LCB}{\mathrm{LCB}}
\newcommand{\fmejor}{f^{\ast}_{n}}
\newcommand{\VAN}{\mathrm{VAN}}
\newcommand{\PSI}{\mathrm{PSI}}
\newcommand{\Emin}{E_{\min}}
\newcommand{\eur}{\text{\,€}}

% --- Partes del examen: capítulos «Parte A», «Parte B», «Parte C» -----------
\renewcommand{\thechapter}{\Alph{chapter}}
\setlocalecaption{spanish}{chapter}{Parte}
\crefname{chapter}{parte}{partes}

% --- Cuestiones del test ----------------------------------------------------
\newlist{cuestiones}{enumerate}{1}
\setlist[cuestiones]{label=\textbf{\arabic*.},ref=\arabic*,leftmargin=2.2em,itemsep=7pt,topsep=6pt}
\newlist{opciones}{enumerate}{1}
\setlist[opciones]{label=(\alph*),leftmargin=2.4em,itemsep=0pt,topsep=2pt}
\newcommand{\pt}[1]{\textit{(#1~p)}\ }

\begin{document}
\frontmatter
% ------------------------------------------------------------------ Portada
\begin{titlepage}
\centering
\vspace*{2.5cm}
{\large\scshape Anatomía de un LLM}\par
\vspace{0.4cm}
{\normalsize Fundamentos matemáticos de los modelos de lenguaje}\par
\vspace{2.2cm}
\rule{\textwidth}{0.8pt}\par\vspace{0.5cm}
{\small\scshape Fase VI \textperiodcentered\ Hito 6 \textperiodcentered\ Examen de hito}\par\vspace{0.3cm}
{\Huge\bfseries Hito 6\par}\vspace{0.3cm}
{\LARGE Diseño y construcción de optimizadores con IA\par}
\vspace{0.5cm}\rule{\textwidth}{0.4pt}\par
\vspace{1.0cm}
{\normalsize Módulos 24 (formulación y resolución), 25 (optimización bayesiana y sustitutos),\\ 26 (control predictivo y aprendizaje por refuerzo), 27 (geometría y 3D)\\ y 28 (búsqueda con LLM, producción y valor)\par}
\vspace{1.4cm}
{\large Examen de hito y proyecto final del curso\par}
\vfill
{\small Edición de octubre de 2026}\par
\end{titlepage}

% =====================================================================
\chapter*{Instrucciones}
\addcontentsline{toc}{chapter}{Instrucciones}
\markboth{Instrucciones}{}
% =====================================================================

Este examen evalúa la fase~VI y cierra el curso. Cubre la formulación y resolución de problemas de optimización (módulo~24), la optimización bayesiana y los modelos sustitutos (módulo~25), el control predictivo y el aprendizaje por refuerzo (módulo~26), la optimización geométrica y 3D (módulo~27) y la búsqueda con modelos de lenguaje, la puesta en producción y la cuantificación del valor (módulo~28). El examen escrito (Partes~A y~B) comprueba el dominio de las herramientas a mano; el proyecto final (Parte~C), que es el núcleo del hito, comprueba que el estudiante sabe diseñar, construir y defender un optimizador con IA completo.

\section*{Condiciones}
\begin{center}
\small
\begin{tabularx}{\textwidth}{@{}lX@{}}
\toprule
Duración & 3 horas (Partes~A y~B). \\
Material permitido & Calculadora científica no programable y una hoja A4 de fórmulas escrita a mano por ambas caras. No se permiten apuntes, libros ni dispositivos conectados. Las tablas y los datos de referencia de estas instrucciones forman parte del examen. \\
Puntuación del examen & 100 puntos: Parte~A, 40 puntos (20 cuestiones de 2 puntos, 4 por módulo); Parte~B, 60 puntos (5 problemas de 12 puntos, uno por módulo). \\
Aprobado del examen & 70 puntos (70\,\%) o más. \\
Proyecto final (Parte~C) & Se califica aparte sobre 100 puntos con la rúbrica de la opción elegida (\cref{cap:proyecto}); dura 8 semanas y se aprueba con 70 puntos o más y todos los criterios de aceptación obligatorios cumplidos. \\
Calificación del hito & Nota del hito $=0{,}4\cdot\text{examen}/10+0{,}6\cdot\text{proyecto}/10$, con al menos un 7 en cada una de las dos partes. \\
Módulos & M24, M25, M26, M27 y M28 (reparto por módulo en el \cref{ap:autoevaluacion}). \\
\bottomrule
\end{tabularx}
\end{center}

\section*{Normas}
\begin{enumerate}[label=\arabic*.]
  \item En la Parte~A hay una sola opción correcta por cuestión. Las respuestas erróneas no restan.
  \item En la Parte~B deben justificarse todos los pasos y escribirse las unidades. Los resultados finales se dan con tres cifras significativas y los intermedios con cuatro. Cuando haga falta $\Phi$ o $\varphi$ fuera de los valores tabulados se interpola linealmente; se admiten diferencias de $\pm0{,}002$ debidas a la interpolación.
  \item Todos los datos de empresas, plantas, escenas y flotas son sintéticos. Las cifras de casos industriales publicados que aparecen en los apuntes no se preguntan de memoria.
  \item Salvo indicación contraria, los problemas de optimización bayesiana son de \emph{minimización} y $\fmejor$ es el menor valor observado, como en M25, tema~4.
\end{enumerate}

\section*{Datos de referencia}

\begin{table}[h]
\centering\footnotesize
\setlength{\tabcolsep}{4pt}
\caption{Densidad $\varphi(z)$ y función de distribución $\Phi(z)$ de la normal estándar. Para $z<0$: $\varphi(-z)=\varphi(z)$ y $\Phi(-z)=1-\Phi(z)$.}\label{tab:normal}
\begin{tabular}{@{}ccc@{\qquad}ccc@{\qquad}ccc@{}}
\toprule
$z$ & $\varphi(z)$ & $\Phi(z)$ & $z$ & $\varphi(z)$ & $\Phi(z)$ & $z$ & $\varphi(z)$ & $\Phi(z)$ \\
\midrule
0{,}00 & 0{,}3989 & 0{,}5000 & 0{,}85 & 0{,}2780 & 0{,}8023 & 1{,}70 & 0{,}0940 & 0{,}9554 \\
0{,}05 & 0{,}3984 & 0{,}5199 & 0{,}90 & 0{,}2661 & 0{,}8159 & 1{,}75 & 0{,}0863 & 0{,}9599 \\
0{,}10 & 0{,}3970 & 0{,}5398 & 0{,}95 & 0{,}2541 & 0{,}8289 & 1{,}80 & 0{,}0790 & 0{,}9641 \\
0{,}15 & 0{,}3945 & 0{,}5596 & 1{,}00 & 0{,}2420 & 0{,}8413 & 1{,}85 & 0{,}0721 & 0{,}9678 \\
0{,}20 & 0{,}3910 & 0{,}5793 & 1{,}05 & 0{,}2299 & 0{,}8531 & 1{,}90 & 0{,}0656 & 0{,}9713 \\
0{,}25 & 0{,}3867 & 0{,}5987 & 1{,}10 & 0{,}2179 & 0{,}8643 & 1{,}95 & 0{,}0596 & 0{,}9744 \\
0{,}30 & 0{,}3814 & 0{,}6179 & 1{,}15 & 0{,}2059 & 0{,}8749 & 2{,}00 & 0{,}0540 & 0{,}9772 \\
0{,}35 & 0{,}3752 & 0{,}6368 & 1{,}20 & 0{,}1942 & 0{,}8849 & 2{,}05 & 0{,}0488 & 0{,}9798 \\
0{,}40 & 0{,}3683 & 0{,}6554 & 1{,}25 & 0{,}1826 & 0{,}8944 & 2{,}10 & 0{,}0440 & 0{,}9821 \\
0{,}45 & 0{,}3605 & 0{,}6736 & 1{,}30 & 0{,}1714 & 0{,}9032 & 2{,}15 & 0{,}0396 & 0{,}9842 \\
0{,}50 & 0{,}3521 & 0{,}6915 & 1{,}35 & 0{,}1604 & 0{,}9115 & 2{,}20 & 0{,}0355 & 0{,}9861 \\
0{,}55 & 0{,}3429 & 0{,}7088 & 1{,}40 & 0{,}1497 & 0{,}9192 & 2{,}25 & 0{,}0317 & 0{,}9878 \\
0{,}60 & 0{,}3332 & 0{,}7257 & 1{,}45 & 0{,}1394 & 0{,}9265 & 2{,}30 & 0{,}0283 & 0{,}9893 \\
0{,}65 & 0{,}3230 & 0{,}7422 & 1{,}50 & 0{,}1295 & 0{,}9332 & 2{,}35 & 0{,}0252 & 0{,}9906 \\
0{,}70 & 0{,}3123 & 0{,}7580 & 1{,}55 & 0{,}1200 & 0{,}9394 & 2{,}40 & 0{,}0224 & 0{,}9918 \\
0{,}75 & 0{,}3011 & 0{,}7734 & 1{,}60 & 0{,}1109 & 0{,}9452 & 2{,}45 & 0{,}0198 & 0{,}9929 \\
0{,}80 & 0{,}2897 & 0{,}7881 & 1{,}65 & 0{,}1023 & 0{,}9505 & 2{,}50 & 0{,}0175 & 0{,}9938 \\
\bottomrule
\end{tabular}
\end{table}

\begin{table}[h]
\centering\small
\caption{Cuantiles y constantes.}\label{tab:cuantiles}
\begin{tabular}{@{}lc@{\qquad}lc@{}}
\toprule
Cantidad & Valor & Cantidad & Valor \\
\midrule
$z_{0{,}75}=\Phi^{-1}(0{,}75)$ & $0{,}6745$ & $t_{15;\,0{,}975}$ (Student, 15 g.\,l.) & $2{,}131$ \\
$z_{0{,}80}$ & $0{,}8416$ & $t_{15;\,0{,}95}$ & $1{,}753$ \\
$z_{0{,}95}$ & $1{,}645$ & $e^{-1/2}$ & $0{,}6065$ \\
$z_{0{,}975}$ & $1{,}960$ & $e^{-2}$ & $0{,}1353$ \\
$\ln 0{,}95$ & $-0{,}05129$ & $e^{-9/2}$ & $0{,}01111$ \\
\bottomrule
\end{tabular}
\end{table}

\section*{Fórmulas}
\[
\EI_n(x)=(\fmejor-\mu)\,\Phi(z)+\sigma\,\varphi(z),\quad z=\frac{\fmejor-\mu}{\sigma};
\]
\[
\mu_n(x)=k_\ast\T K_y^{-1}y,\qquad \sigma_n^2(x)=k(x,x)-k_\ast\T K_y^{-1}k_\ast,\qquad K_y=K+\sigma_\varepsilon^{2}I ;
\]
\[
\VAN=-I+b\,\frac{1-(1+\rho)^{-N}}{\rho};\qquad
n\approx\Bigl(\frac{(z_{1-\alpha/2}+z_{1-\beta})\,\sigma_D}{\delta}\Bigr)^{2};\qquad
\bar d\pm t_{n-1,\,1-\alpha/2}\,\frac{s_d}{\sqrt n}.
\]


\tableofcontents

\mainmatter
% =====================================================================
\chapter{Cuestiones}
\label{cap:test}
% =====================================================================

Veinte cuestiones de opción múltiple, de 2 puntos cada una (40 puntos en total), cuatro por módulo. Hay una sola respuesta correcta en cada cuestión y las respuestas erróneas no restan. La clave y su justificación están en el \cref{ap:soluciones}.

\section*{Módulo 24. Formulación y resolución}
\phantomsection\addcontentsline{toc}{section}{Módulo 24. Formulación y resolución}

\begin{cuestiones}
\item \pt{2} Una línea de envasado puede producir como máximo 120~t por semana. Con inventario final nulo, la demanda pendiente desde la semana~$t$ hasta el final del horizonte es de 90~t. En la restricción $x_t\le M\,y_t$, con $y_t\in\{0,1\}$ la variable de preparación, la menor $M$ válida es:
\begin{opciones}
  \item $120$~t;
  \item $210$~t;
  \item $90$~t;
  \item $10^{6}$~t, para no cortar ninguna solución.
\end{opciones}

\item \pt{2} En el óptimo de un programa lineal que maximiza el margen semanal, a la restricción de horas de granuladora le sobran 12~h. Su precio sombra:
\begin{opciones}
  \item es cero;
  \item es el margen óptimo dividido entre 12~h;
  \item coincide con el margen por hora del producto que más granuladora consume;
  \item no puede conocerse sin resolver el dual completo.
\end{opciones}

\item \pt{2} La demanda semanal de un producto perecedero es $\N(200,\,40^{2})$~hL. Cada hectolitro que falta cuesta $c_u=30$~€ y cada hectolitro que sobra, $c_o=10$~€. La cantidad óptima es, aproximadamente:
\begin{opciones}
  \item $200$~hL;
  \item $173$~hL;
  \item $234$~hL;
  \item $227$~hL.
\end{opciones}

\item \pt{2} Sean dos objetivos que se minimizan y pesos $w_1,w_2>0$. Si $x^{\ast}$ minimiza $w_1f_1+w_2f_2$ en el conjunto factible:
\begin{opciones}
  \item $x^{\ast}$ puede estar dominado por otro punto factible;
  \item $x^{\ast}$ es Pareto-óptimo, aunque puede haber puntos del frente que ninguna suma ponderada alcance;
  \item variando los pesos se obtiene siempre el frente de Pareto completo;
  \item $x^{\ast}$ es Pareto-óptimo solo si $f_1$ y $f_2$ son convexas.
\end{opciones}
\end{cuestiones}

\section*{Módulo 25. Optimización bayesiana y modelos sustitutos}
\phantomsection\addcontentsline{toc}{section}{Módulo 25. Optimización bayesiana y modelos sustitutos}

\begin{cuestiones}[resume]
\item \pt{2} En un problema de minimización, el mejor valor observado es $\fmejor=10$ y la posterior en un candidato es $f(x)\sim\N(10,\,2^{2})$. Su mejora esperada vale:
\begin{opciones}
  \item $0$, porque la media no mejora $\fmejor$;
  \item $0{,}399$;
  \item $0{,}798$;
  \item $2$.
\end{opciones}

\item \pt{2} Un proceso gaussiano de media nula con núcleo exponencial cuadrático, $\sigma_f=\ell=1$, se condiciona a una sola observación sin ruido en $x=0$. Varianza posterior de $f(1)$:
\begin{opciones}
  \item $0{,}632$;
  \item $0{,}368$;
  \item $1$, porque la observación no informa sobre otros puntos;
  \item $0{,}135$.
\end{opciones}

\item \pt{2} En una búsqueda aleatoria de hiperparámetros se quiere que, con probabilidad $0{,}95$, al menos una configuración caiga en la región del mejor $5\,\%$. Número mínimo de configuraciones independientes:
\begin{opciones}
  \item $20$;
  \item $19$;
  \item $100$;
  \item $59$.
\end{opciones}

\item \pt{2} Una biblioteca devuelve la desviación típica de una nueva \emph{observación}, que es $\sqrt{\sigma_n^{2}(x)+\sigma_\varepsilon^{2}}$, y por error se usa en la mejora esperada de un problema con ruido. Entonces, en un punto ya ensayado:
\begin{opciones}
  \item la mejora esperada se anula, como debe ser;
  \item la mejora esperada no baja del nivel que fija el ruido y el algoritmo tiende a repetir ensayos casi iguales;
  \item la mejora esperada se vuelve negativa;
  \item no hay efecto apreciable si $\sigma_\varepsilon$ es menor que $\sigma_f$.
\end{opciones}
\end{cuestiones}

\section*{Módulo 26. Control predictivo y aprendizaje por refuerzo}
\phantomsection\addcontentsline{toc}{section}{Módulo 26. Control predictivo y aprendizaje por refuerzo}

\begin{cuestiones}[resume]
\item \pt{2} Un modelo térmico continuo tiene un único autovalor $\lambda=-5$~h$^{-1}$. El modelo discreto obtenido con Euler explícito es estable si y solo si el periodo de muestreo cumple:
\begin{opciones}
  \item $\Delta t<0{,}4$~h;
  \item $\Delta t<0{,}2$~h;
  \item cualquier $\Delta t>0$, porque el sistema continuo es estable;
  \item $\Delta t<0{,}5$~h.
\end{opciones}

\item \pt{2} Sea $x_{k+1}=x_k+u_k$ con coste $x_0^{2}+u_0^{2}+x_1^{2}$ (horizonte $N=1$, $Q=R=P_f=1$) y $x_0=2$. La decisión óptima es:
\begin{opciones}
  \item $u_0=-2$;
  \item $u_0=-0{,}5$;
  \item $u_0=-1$;
  \item $u_0=-1{,}2$.
\end{opciones}

\item \pt{2} En un MPC convexo, la restricción de confort tiene en el óptimo un multiplicador de Karush--Kuhn--Tucker de 3~€ por K·h. Se ablanda con una holgura penalizada con $\rho\sum_j s_j$. La solución blanda coincide con la dura, cuando esta es factible, si:
\begin{opciones}
  \item $\rho>0$;
  \item $\rho\ge 1/3$;
  \item $\rho\le 3$;
  \item $\rho\ge 3$.
\end{opciones}

\item \pt{2} Un episodio de climatización se corta a las 72~h por un límite de tiempo (\texttt{truncated}) sin que el proceso haya terminado. El objetivo del crítico en la última transición $(s,a,r,s')$ debe ser:
\begin{opciones}
  \item $y=r$;
  \item $y=r+\gamma\,Q_{\phi'}(s',\tilde a')$;
  \item $y=\gamma\,Q_{\phi'}(s',\tilde a')$;
  \item $y=r-\gamma\,Q_{\phi'}(s',\tilde a')$.
\end{opciones}
\end{cuestiones}

\section*{Módulo 27. Optimización geométrica y 3D}
\phantomsection\addcontentsline{toc}{section}{Módulo 27. Optimización geométrica y 3D}

\begin{cuestiones}[resume]
\item \pt{2} Con la interpolación SIMP, $p=3$ y $\Emin$ despreciable, un elemento de densidad $x_e=0{,}6$ aporta una rigidez relativa a la del sólido de:
\begin{opciones}
  \item $0{,}6$;
  \item $0{,}36$;
  \item $0{,}216$;
  \item $0{,}4$.
\end{opciones}

\item \pt{2} En el problema de flexibilidad mínima, $c=F\T U$ con $K(x)U=F$ y $N$ densidades. Además de la resolución de $KU=F$, el gradiente $\nabla_xc$ completo necesita:
\begin{opciones}
  \item ninguna resolución adicional, porque el adjunto es $\lambda=U$;
  \item una resolución adjunta con otro término independiente;
  \item $N$ resoluciones, una por densidad;
  \item $2N$ resoluciones por diferencias centradas.
\end{opciones}

\item \pt{2} Una cámara con campo de visión vertical de $60^{\circ}$ produce imágenes de $1080$ píxeles de alto. Un nivel de detalle con error geométrico de $0{,}02$~m puede usarse sin superar $1$~píxel de error a partir de una distancia de:
\begin{opciones}
  \item $10{,}8$~m;
  \item $37{,}4$~m;
  \item $32{,}4$~m;
  \item $18{,}7$~m.
\end{opciones}

\item \pt{2} Una malla triangular cerrada con topología de esfera ($\chi=2$) tiene 1000 vértices. Su número de caras es:
\begin{opciones}
  \item $1996$;
  \item $2000$;
  \item $2994$;
  \item $998$.
\end{opciones}
\end{cuestiones}

\section*{Módulo 28. Búsqueda con LLM, producción y valor}
\phantomsection\addcontentsline{toc}{section}{Módulo 28. Búsqueda con LLM, producción y valor}

\begin{cuestiones}[resume]
\item \pt{2} Cincuenta instrucciones de un clasificador tienen la misma exactitud verdadera y se evalúan con errores $\N(0,\,0{,}04^{2})$. La cota de la maldición del optimizador para el sesgo de la instrucción ganadora es:
\begin{opciones}
  \item $0{,}040$;
  \item $0{,}112$;
  \item $0{,}280$;
  \item $0{,}0057$.
\end{opciones}

\item \pt{2} Una variable binaria de un sensor tenía proporciones $p=(0{,}5;\ 0{,}5)$ en entrenamiento y $q=(0{,}6;\ 0{,}4)$ en la última semana. Su índice de estabilidad poblacional es:
\begin{opciones}
  \item $0{,}0203$;
  \item $0{,}10$;
  \item $0$, porque las medias coinciden;
  \item $0{,}0405$.
\end{opciones}

\item \pt{2} Un halving sucesivo con $\eta=3$ y $K=3$ escalones ($r_{\max}=27\,r_{\min}$, $N$ múltiplo de 27 y coste proporcional al presupuesto) cuesta, respecto a evaluar las $N$ configuraciones con el presupuesto completo, una fracción de:
\begin{opciones}
  \item $1/27$;
  \item $1/3$;
  \item $4/27$;
  \item $1/4$.
\end{opciones}

\item \pt{2} Un optimizador cuesta $I=50\,000$~€ y aporta un flujo neto constante de $20\,000$~€ al final de cada uno de 4~años. Con un tipo de descuento del $10\,\%$, su valor actual neto es, aproximadamente:
\begin{opciones}
  \item $30\,000$~€;
  \item $13\,400$~€;
  \item $63\,400$~€;
  \item $-6\,600$~€.
\end{opciones}
\end{cuestiones}

% =====================================================================
\chapter{Problemas}
\label{cap:problemas}
% =====================================================================

Cinco problemas de 12 puntos cada uno (60 puntos en total), uno por módulo. La puntuación de cada apartado figura entre paréntesis. Los problemas marcados con \emph{Fuente} están adaptados de exámenes, hojas de problemas o trabajos de curso de universidades; el enunciado es de redacción propia y los datos se han cambiado para que puedan resolverse con calculadora. Tres de ellos (B.1, B.3 y B.5) proceden de exámenes o de colecciones de preparación de exámenes. Las soluciones completas están en el \cref{ap:soluciones}.

% ---------------------------------------------------------------------
\section*{Problema B.1. Hornos, costes fijos y una disyunción (M24)}
\phantomsection\addcontentsline{toc}{section}{Problema B.1. Hornos, costes fijos y una disyunción}

\begin{ejercicio}{2}{eje:milp}
Una planta de tratamiento térmico debe procesar exactamente $D=100$~t en un turno con tres hornos. Encender el horno $i$ cuesta $f_i$ (precalentamiento y energía fija); procesar cada tonelada cuesta $c_i$ en energía; si se enciende, por estabilidad del proceso debe tratar al menos $L_i$ toneladas, y nunca más de su capacidad $U_i$:
\begin{center}\small
\begin{tabular}{@{}lcccc@{}}
\toprule
Horno & $f_i$ (€) & $c_i$ (€/t) & $L_i$ (t) & $U_i$ (t) \\
\midrule
1 & 400 & 12 & 30 & 80 \\
2 & 150 & 16 & 20 & 50 \\
3 & 250 & 14 & 25 & 60 \\
\bottomrule
\end{tabular}
\end{center}
\begin{enumerate}[label=(\alph*)]
  \item \pt{3} Formule el problema de coste mínimo como un programa lineal entero mixto con variables $x_i\ge 0$ (toneladas) e $y_i\in\{0,1\}$ (encendido), usando restricciones de tipo big-M. Indique qué representa cada restricción.
  \item \pt{2} Determine la menor constante $M_i$ válida de cada restricción $x_i\le M_iy_i$ y justifique que no elimina ninguna solución factible.
  \item \pt{3} Resuelva el problema enumerando los patrones de encendido $y\in\{0,1\}^{3}$; para cada patrón factible, resuelva a mano el programa lineal en $x$. Dé el plan óptimo y su coste.
  \item \pt{2} Calcule el valor de la relajación lineal ($y_i\in[0,1]$) con las $M_i$ de~(b) y el \emph{gap} de integralidad. Repita el cálculo si se escribe la capacidad como restricción aparte, $x_i\le U_i$, y el enlace con $M=10^{4}$ para los tres hornos. ¿Qué concluye?
  \item \pt{2} La acometida eléctrica impone que, si el horno~2 está encendido, el horno~1 no trate más de 60~t. Escriba la restricción lineal con la menor constante válida y obtenga el nuevo óptimo.
\end{enumerate}
\fuente{apartados~(a), (b) y~(e) adaptados de MIT 15.053, primavera de 2013, colección de problemas de preparación del segundo examen parcial (\emph{Midterm Exam~2: Practice Problem Set}), problema~2 (apertura de cajas de cobro con coste fijo y restricciones de enlace) y problema~4(a) (restricción disyuntiva «al menos una de dos»)~\cite{mit15053mid}; apartado~(d) inspirado en la hoja de problemas~5, problema~4, del mismo curso (elegir la menor $M$ válida)~\cite{mit15053ps}.}
\end{ejercicio}

% ---------------------------------------------------------------------
\section*{Problema B.2. Un paso de optimización bayesiana a mano (M25)}
\phantomsection\addcontentsline{toc}{section}{Problema B.2. Un paso de optimización bayesiana a mano}

\begin{ejercicio}{2}{eje:bo}
Se ajusta la temperatura de curado de un adhesivo (variable codificada $x$) para minimizar el tiempo de ciclo tipificado. Se modeliza $f\sim\mathcal{GP}(0,k)$ con núcleo exponencial cuadrático $k(x,x')=\exp\bigl(-\tfrac12(x-x')^{2}/\ell^{2}\bigr)$, $\sigma_f=1$ y $\ell=1$. Se han hecho dos ensayos sin ruido: $f(0)=0{,}5$ y $f(2)=-0{,}8$.
\begin{enumerate}[label=(\alph*)]
  \item \pt{3} Escriba $K$, calcule $K^{-1}$ y $\alpha=K^{-1}y$, y obtenga la media y la varianza posteriores en $x=1$.
  \item \pt{2} Obtenga la media y la varianza posteriores en $x=3$.
  \item \pt{3} Calcule la probabilidad de mejora y la mejora esperada en $x=1$ y en $x=3$, y la cota inferior de confianza $\mu-2\sigma$ en ambos. ¿Qué punto se ensaya a continuación con cada criterio? Explique el resultado en términos de explotación y exploración.
  \item \pt{2} Si los ensayos tuvieran ruido gaussiano de varianza $\sigma_\varepsilon^{2}=0{,}1$, calcule la varianza posterior de $f(2)$ y la de una nueva observación $y(2)$. ¿Cuál debe usarse en la mejora esperada y por qué?
  \item \pt{2} Sin ruido, se cambia la longitud de escala a $\ell=0{,}3$. Calcule aproximadamente $\mu(1)$, $\sigma^{2}(1)$ y $\EI(1)$, y explique cómo afecta la longitud de escala a la decisión.
\end{enumerate}
\fuente{adaptado del trabajo de curso~1, «Gaussian processes», de Cambridge 4F13 \emph{Probabilistic Machine Learning} (Michaelmas de 2024), que pide predecir con un proceso gaussiano, interpretar la incertidumbre y estudiar el efecto de los hiperparámetros~\cite{cam4f13}, y del ejemplo de la mejora esperada de Jones, Schonlau y Welch~\cite{jones}.}
\end{ejercicio}

% ---------------------------------------------------------------------
\section*{Problema B.3. Riccati, MPC de un paso y Bellman con coste lineal (M26)}
\phantomsection\addcontentsline{toc}{section}{Problema B.3. Riccati, MPC de un paso y Bellman con coste lineal}

\begin{ejercicio}{3}{eje:mpc}
La desviación de temperatura de un depósito de fluido térmico respecto a su consigna, $x_k$ (K), evoluciona con $x_{k+1}=a\,x_k+b\,u_k$, donde $u_k$ es la potencia de calefacción o refrigeración (kW, en variables de desviación).
\begin{enumerate}[label=(\alph*)]
  \item \pt{3} Con $a=0{,}9$, $b=0{,}5$ K/kW, coste $\sum_{j=0}^{1}(Qx_j^{2}+Ru_j^{2})+P_fx_2^{2}$, $Q=1$, $R=0{,}5$, $P_f=1$ y $x_0=2$~K, aplique la recursión de Riccati para obtener $K_1$, $P_1$, $K_0$ y $P_0$, la secuencia óptima sin restricciones y su coste. Compruebe que coincide con $P_0x_0^{2}$.
  \item \pt{3} El actuador está limitado a $\abs{u_j}\le 1$~kW. Escriba el coste como $J(U)=U\T HU+2x_0F\T U+\text{cte}$, resuelva el MPC de horizonte 2, indique qué decisión se aplica en el instante actual y calcule el multiplicador de la restricción activa.
  \item \pt{4} En otra operación se usa el coste descontado con término lineal
  \[
  J=\sum_{t=0}^{\infty}\gamma^{t}\bigl(q\,x_t+r\,u_t^{2}\bigr),\qquad q=2,\ r=1,\ \gamma=0{,}95,
  \]
  sin restricciones y con los mismos $a$ y $b$. Plantee la ecuación de Bellman para $V(z)=\min J$ con $x_0=z$, compruebe que una función afín $V(z)=pz+s$ la satisface y calcule $p$, $s$ y la acción óptima.
  \item \pt{2} Se añade a~(c) la restricción $\abs{u_t}\le 2$~kW. Explique por qué $V$ sigue siendo afín con la misma pendiente $p$, y calcule la nueva acción óptima y el nuevo término $s$.
\end{enumerate}
\fuente{apartado~(c) adaptado de Stanford EE363 \emph{Linear Dynamical Systems} (S.~Boyd), examen final de invierno de 2006, problema~8, «LQR with linear state cost», escalarizado~\cite{ee363final}; apartados~(a) y~(b) siguen la clase~16 de EE363 y el capítulo~1 de Rawlings, Mayne y Diehl~\cite{ee363,rmd}.}
\end{ejercicio}

% ---------------------------------------------------------------------
\section*{Problema B.4. Sensibilidad adjunta y SIMP en una barra (M27)}
\phantomsection\addcontentsline{toc}{section}{Problema B.4. Sensibilidad adjunta y SIMP en una barra}

\begin{ejercicio}{3}{eje:simp}
Una barra empotrada en su extremo izquierdo (nodo~0) se discretiza en tres elementos en serie con nodos libres~1, 2 y~3. El elemento $e$ tiene rigidez axial $k_e=x_e^{\,p}k_0$, con $k_0=1000$~N/mm, $p=3$ y densidades $x=(1{,}0;\ 0{,}8;\ 0{,}5)$ ($\Emin$ despreciable). Actúan cargas de 100~N en los nodos~2 y~3, es decir, $F=(0,100,100)$~N.
\begin{enumerate}[label=(\alph*)]
  \item \pt{2} Ensamble la matriz de rigidez $K\in\R^{3\times 3}$ (tridiagonal), obtenga los desplazamientos $U$ y la flexibilidad $c=F\T U$.
  \item \pt{3} Calcule $\partial c/\partial x_e$ para los tres elementos con el método adjunto, justificando que el problema es autoadjunto. Compruebe uno de los valores derivando directamente la expresión de $c$.
  \item \pt{3} Un sensor mide el desplazamiento del nodo~2 y se quiere controlar $J=u_2=e_2\T U$. Plantee y resuelva el problema adjunto y calcule $\partial J/\partial x_e$. Interprete el valor obtenido para el elemento~3.
  \item \pt{2} Se aplica una actualización por criterios de optimalidad con $\eta=1/2$, límite de movimiento $m=0{,}2$ y volumen fijo $\sum_ex_e=2{,}3$. Calcule $x^{\text{nuevo}}$ con $\Lambda=300$, compruebe si cumple el volumen y diga hacia dónde debe moverse $\Lambda$. Verifique que $\Lambda=270$ cumple la restricción.
  \item \pt{2} Con $p=1$ el problema es convexo. Halle con las condiciones de Karush--Kuhn--Tucker las densidades que minimizan $c$ con $\sum_ex_e=2{,}3$ (ignore las cotas $x_e\le 1$ y compruebe después que se cumplen) y compare la flexibilidad con la del diseño inicial evaluado con $p=1$.
\end{enumerate}
\fuente{apartados~(b) y~(c) adaptados de MIT 18.S096 \emph{Matrix Calculus for Machine Learning and Beyond} (IAP 2023), hoja de problemas~2, problema~1 (gradiente de una función de la solución de un sistema tridiagonal con una sola resolución adjunta)~\cite{mit18s096ps2}; apartados~(d) y~(e) siguen el capítulo~1 de Bendsøe y Sigmund~\cite{bs}.}
\end{ejercicio}

% ---------------------------------------------------------------------
\section*{Problema B.5. ¿Compensa el optimizador? Intervalos y valor actual (M28)}
\phantomsection\addcontentsline{toc}{section}{Problema B.5. ¿Compensa el optimizador? Intervalos y valor actual}

\begin{ejercicio}{2}{eje:valor}
Una planta ha probado durante $n=16$ semanas un optimizador de planificación energética. Cada semana se calcula, con el gemelo digital calibrado y los mismos pedidos, precios y averías, la diferencia de coste $d_i=$ (coste con el plan actual) $-$ (coste con el optimizador). Se obtiene $\bar d=1200$~€/semana y $s_d=1600$~€/semana.
\begin{enumerate}[label=(\alph*)]
  \item \pt{3} Calcule el intervalo de confianza del 95\,\% para la mejora media semanal y expréselo en €/año (52 semanas). Enuncie las hipótesis que lo justifican.
  \item \pt{2} Por separado, los costes semanales de cada política tienen desviación típica de 3000~€ y correlación $0{,}86$. Calcule la desviación típica de la diferencia pareada y la que tendría la diferencia con semanas independientes. ¿Cuántas veces más semanas harían falta sin emparejar para la misma precisión?
  \item \pt{2} Con $B=1000$ remuestras \emph{bootstrap} de las 16 diferencias, los percentiles 5 y 95 de las medias remuestreadas son 560 y 1890~€/semana. Dé los intervalos \emph{bootstrap} percentil y básico del 90\,\%. ¿Por qué pueden no ser simétricos respecto a $\bar d$?
  \item \pt{3} La inversión es $I=60\,000$~€ y la operación cuesta $12\,000$~€/año; el horizonte es de 5 años con $\rho=8\,\%$. Calcule el VAN con la mejora media y con el extremo inferior del intervalo de~(a), y la mejora semanal mínima que hace $\VAN=0$. ¿Puede afirmarse con un 95\,\% de confianza que el proyecto es rentable?
  \item \pt{2} ¿Cuántas semanas pareadas harían falta para detectar, con $\alpha=0{,}05$ bilateral y potencia $0{,}8$, una mejora media igual a la de rentabilidad nula de~(d), suponiendo $\sigma_D=1600$~€?
\end{enumerate}
\fuente{adaptado de MIT 18.05 \emph{Introduction to Probability and Statistics} (Orloff y Bloom), examen final de primavera de 2022, problemas~9 (contraste $t$ e intervalo de confianza), 11 (potencia de un contraste) y~13 (intervalo \emph{bootstrap} percentil)~\cite{mit1805final}; apartado~(d), de la sesión «Present Value Relations» de MIT 15.401~\cite{mit15401}.}
\end{ejercicio}

% =====================================================================
\chapter{Proyecto final: diseñe y construya un optimizador con IA}
\label{cap:proyecto}
% =====================================================================

El proyecto final cierra el curso. El estudiante diseña, construye, evalúa y defiende un optimizador con IA completo para uno de tres ámbitos: una planta industrial (opción~a), un \emph{pipeline} de software 3D (opción~b) o las operaciones de una flota (opción~c). En los tres hay que recorrer los doce pasos de M28, §6.1: decidir qué se optimiza y contra qué se compara, generar datos sintéticos realistas y un gemelo digital, formular el problema, combinar un resolvedor con un componente aprendido, proteger las restricciones duras y demostrar con intervalos de confianza cuánto mejora el sistema y si compensa. La \cref{fig:arquitectura} resume la arquitectura que se exige en las tres opciones; es la de M24, §1.2, completada con la evaluación de M28, temas~4 y~5.

\begin{figure}[h]
\centering
\begin{tikzpicture}[font=\footnotesize,>={Stealth[length=2mm]},node distance=4mm and 4mm,
  caja/.style={draw,thick,rounded corners=2pt,align=center,inner sep=2pt,minimum height=10mm,text width=2.25cm}]
  \node[caja] (gen) {generador de datos\\(semilla fija)};
  \node[caja,right=of gen] (gem) {gemelo digital\\(simulador)};
  \node[caja,right=of gem] (pred) {modelo predictivo\\o sustituto};
  \node[caja,right=of pred] (opt) {optimizador\\(PLEM, MPC, BO\dots)};
  \node[caja,below=12mm of opt] (seg) {salvaguardas\\y respaldo};
  \node[caja,left=of seg] (dec) {decisión\\aplicada};
  \node[caja,left=of dec] (eva) {evaluación\\con CRN\\e intervalos};
  \node[caja,left=of eva] (mon) {monitorización\\y valor (VAN)};
  \draw[->,thick] (gen)--(gem); \draw[->,thick] (gem)--(pred); \draw[->,thick] (pred)--(opt);
  \draw[->,thick] (opt)--(seg); \draw[->,thick] (seg)--(dec); \draw[->,thick] (dec)--(eva); \draw[->,thick] (eva)--(mon);
  \draw[->,thick,dashed] (dec.north) -- ++(0,6mm) -| node[pos=0.13,above=-0.5pt,font=\scriptsize] {se aplica en el gemelo} (gem.south);
\end{tikzpicture}
\caption{Arquitectura mínima exigida en las tres opciones del proyecto. Las flechas continuas forman el ciclo de decisión; la discontinua indica que las decisiones se evalúan en el gemelo, nunca solo con el modelo que las propuso.}
\label{fig:arquitectura}
\end{figure}

% =====================================================================
\section{Organización, reglas y entregables comunes}
\label{sec:proy-comun}

\paragraph{Formato.} Trabajo individual o en parejas, con el reparto de tareas declarado en el documento de diseño. Duración: 8 semanas. Se elige una sola opción en la semana~1 y no se cambia después de la semana~2. Cada opción describe el problema, los datos sintéticos que hay que generar, la formulación que se pide, los métodos y tecnologías sugeridos, los hitos semanales, los entregables específicos, los criterios de aceptación y la rúbrica de 100 puntos.

\paragraph{Reglas de evaluación honesta.} Son obligatorias en las tres opciones y se comprueban en la corrección (\cref{ap:correccion}).
\begin{enumerate}[label=R\arabic*.]
  \item \emph{Datos sintéticos con semilla.} Todos los datos se generan con un programa propio y semillas documentadas; no se usan datos reales confidenciales. Los parámetros del generador que se dan en cada opción son el escenario nominal; se pueden enriquecer, pero no simplificar.
  \item \emph{Partición temporal congelada.} La partición en entrenamiento, validación y prueba se fija en la semana~1, y el conjunto de prueba se genera, se guarda con su huella (\emph{hash}) y no se usa para ninguna elección hasta la semana~7 (M08, §1.4; M28, §2.2, corolario~2.6).
  \item \emph{Tres referencias.} Toda mejora se mide frente a la \emph{línea base} (la práctica actual que se describe en cada opción), frente a la \emph{mejor alternativa sencilla} sin IA y frente al \emph{techo} con información perfecta (M28, §5.1 y proposición~6.1).
  \item \emph{Comparación pareada.} Todas las políticas se evalúan en el gemelo con números aleatorios comunes (M26, §5.1), en las mismas semanas o trayectorias, con al menos 5 semillas.
  \item \emph{Intervalos de la mejora.} Se informa la mejora media con intervalo $t$ pareado del 95\,\% y con intervalo \emph{bootstrap} (por bloques si hay autocorrelación) (M28, §5.2), nunca solo con una media.
  \item \emph{Salvaguardas independientes.} Las restricciones duras se comprueban fuera del optimizador antes de aplicar cada decisión, y existe una política de respaldo (M26, §6.5; M28, §5.5).
  \item \emph{Uso de modelos de lenguaje.} Se permite usar asistentes de código y búsquedas guiadas por LLM (M28, tema~1); se declara en el informe qué se generó así y cómo se verificó. El código generado se ejecuta aislado.
\end{enumerate}

\paragraph{Entregables comunes.} Las tres opciones entregan:
\begin{enumerate}[label=E\arabic*.]
  \item \emph{Documento de diseño} con la plantilla del \cref{ap:plantilla}: versión~0 en la semana~2 (secciones 1 a~6), versión~1 en la semana~5 y versión final en la semana~8; entre 10 y 15 páginas más la ficha resumen.
  \item \emph{Código} en un repositorio con estructura \texttt{datos/}, \texttt{src/}, \texttt{experimentos/}, \texttt{pruebas/} y \texttt{informe/}; un \texttt{README} con instrucciones; una única orden (por ejemplo, \texttt{python -m proyecto reproducir --semilla 0}) que regenera datos, modelos, decisiones, tablas y figuras en menos de 2~h en la CPU de un portátil (o menos de 20~min en un modo reducido documentado); pruebas automáticas (\texttt{pytest}) de las restricciones duras, del generador y del cálculo de costes; registro de experimentos (MLflow o equivalente).
  \item \emph{Informe técnico final} de 12 a 20 páginas: resumen, resultados frente a las tres referencias con intervalos de confianza, ablaciones (qué aporta cada componente), pruebas de estrés, valor actual neto con incertidumbre, AMFE, monitorización, limitaciones y trabajo futuro.
  \item \emph{Demostración} de 15~minutos más 10 de preguntas: un ciclo de decisión completo en directo o grabado, la activación de una salvaguarda ante una entrada anómala inyectada y el panel de seguimiento. Las preguntas pueden pedir cambiar un parámetro del generador y volver a ejecutar.
\end{enumerate}

\paragraph{Calendario común.} Las tres opciones siguen el mismo esqueleto de ocho semanas (\cref{tab:calendario}); cada opción lo concreta en su tabla de hitos.

\begin{table}[H]
\centering\small
\caption{Esqueleto común de las ocho semanas.}\label{tab:calendario}
\begin{tabularx}{\textwidth}{@{}c>{\raggedright\arraybackslash}X>{\raggedright\arraybackslash}p{4.6cm}@{}}
\toprule
Semana & Objetivo & Evidencia en el repositorio \\
\midrule
1 & Concepto de operación, requisitos, generador de datos v1, línea base en el gemelo, partición congelada & generador, \emph{hash} de la prueba, coste de la línea base \\
2 & Gemelo validado, techo de valor y mejor alternativa sencilla & documento de diseño v0 \\
3 & Formulación principal resuelta con información perfecta & modelo, tiempos de cálculo, \emph{gap} \\
4 & Componente predictivo con validación temporal e integración & métricas de predicción y de decisión \\
5 & Segundo nivel de decisión (control, selección en línea o rutas) y salvaguardas & documento de diseño v1 \\
6 & Ajuste (optimización bayesiana o multifidelidad) y pruebas de estrés & informe de estrés \\
7 & Evaluación final en prueba con CRN, intervalos, VAN, AMFE y monitorización & tablas finales \\
8 & Informe, API o herramienta, panel, demostración y reproducibilidad & entrega completa \\
\bottomrule
\end{tabularx}
\end{table}

% =====================================================================
\section{Opción (a): planta industrial}
\label{sec:proy-a}

\subsection*{Planificación y control energético de una línea de producción}

\paragraph{Contexto y decisiones.} Una línea fabrica piezas de aluminio en cuatro familias, P1 a P4. Cada lote pasa por un horno de tratamiento térmico (una carga a la vez) y después por uno de dos centros de mecanizado. La energía eléctrica se compra a precio horario y la potencia contratada es de 1000~kW. El optimizador toma dos decisiones encadenadas: (i) cada día a las 14:00, con los precios del día siguiente ya publicados, el \emph{plan horario} de los siete días siguientes (qué lote entra en el horno y en cada máquina en cada hora, cuándo se enciende o se mantiene caliente el horno), en horizonte deslizante, y (ii) cada 5~minutos, la \emph{potencia del horno} que sigue la consigna de temperatura del plan con el menor coste y sin violar la potencia contratada ni la banda de temperatura del tratamiento.

\paragraph{Datos sintéticos que hay que generar.} Un año ficticio de 52 semanas con resolución horaria (planificación) y de 5~minutos (horno).
\begin{enumerate}[label=(\roman*)]
  \item \emph{Precio de la electricidad} (€/MWh): $p_t=\max\bigl(0,\ \mu_{d(t)}+g(h_t,d(t))+\varepsilon_t\bigr)$, con nivel diario $\mu_d=75+15\cos\bigl(2\pi(d-15)/365\bigr)+\eta_d$, $\eta_d=0{,}8\,\eta_{d-1}+\N(0,8^{2})$; perfil horario $g$ con picos de $+20$~€/MWh de 8 a 10~h y de 19 a 22~h y un valle solar de $-25$~€/MWh de 12 a 16~h entre abril y septiembre; ruido $\varepsilon_t\sim\N(0,6^{2})$ y, con probabilidad $0{,}5\,\%$ por hora, un pico que multiplica el precio por~3.
  \item \emph{Temperatura ambiente} (°C): $T^{\mathrm{amb}}_t=16+9\sin\bigl(2\pi(d-105)/365\bigr)+5\sin\bigl(2\pi(h-9)/24\bigr)+\N(0,1{,}5^{2})$.
  \item \emph{Pedidos}: llegadas diarias de Poisson con medias semanales de 10, 7, 5 y 3 lotes de P1 a P4 (factor $1{,}15$ de marzo a junio), fecha de entrega uniforme entre 3 y 10 días laborables después de la llegada. La utilización media resultante del horno debe quedar entre el 70 y el 85\,\%; se ajustan las medias si no es así.
  \item \emph{Horno}: modelo térmico $C\,\dot T=P-(T-T^{\mathrm{amb}})/R-P^{\mathrm{carga}}$, con $R=1$~K/kW, $C=3$~kWh/K ($\tau=3$~h), $P\in[0,800]$~kW, y una carga fría que absorbe 150~kWh durante la primera hora. Tratamientos de 2, 3, 4 y 2~h a 480, 530, 500 y 450~°C, con tolerancia de $\pm5$~°C (restricción dura de calidad). Medida con ruido de 2~°C; un disparo del horno al mes, de 1 a 4~h.
  \item \emph{Mecanizado}: dos centros de 150~kW cada uno en funcionamiento, tiempos de 1 a 3~h por lote según la familia, solo en turnos de lunes a viernes de 6 a 22~h; el horno puede trabajar de noche, pero se carga y descarga solo en turno.
  \item \emph{Costes}: energía $p_t\,E_t$; arranque del horno en frío (desde menos de 200~°C), 120~€; retraso, 60~€ por lote y día; almacenamiento intermedio, 3~€ por lote y día; superar los 1000~kW está prohibido (restricción dura).
  \item \emph{Partición}: semanas 1 a~30 para entrenamiento, 31 a~40 para validación y 41 a~52 para prueba.
\end{enumerate}

\paragraph{Referencias de comparación.} \emph{Línea base}: lotes por fecha de entrega más próxima (EDD) en cuanto hay recurso, y el horno con un PI ajustado por SIMC (M26, §1.5) a consigna fija en turno y apagado fuera de turno. \emph{Mejor alternativa sencilla}: EDD con precalentamiento en las dos horas más baratas antes del turno y un PI con antisaturación. \emph{Techo}: el plan óptimo con precios y pedidos reales conocidos de antemano.

\paragraph{Formulación que se pide.}
\begin{enumerate}[label=F\arabic*.]
  \item \emph{Plan horario} como PLEM en horizonte deslizante de 168~h: binarias de inicio de cada lote en el horno y en cada máquina, estado del horno (encendido, mantenimiento en caliente, apagado) y arranques; restricciones de precedencia horno–mecanizado, de un lote a la vez por recurso, de turnos, de balance de potencia $P_t^{\mathrm{horno}}+\sum_kP_{kt}^{\mathrm{mec}}\le 1000$~kW y de retraso $T_l\ge C_l-d_l$; objetivo de energía prevista, retrasos, arranques y almacenamiento. Todas las constantes big-M deben ser las mínimas válidas (proposición~2.11 de M24) y justificarse. Se informa del tamaño, del \emph{gap} y del tiempo de cálculo, y del precio sombra de la potencia contratada con las binarias fijas (M24, §2.2).
  \item \emph{Incertidumbre del precio}: a partir de D+2, el plan usa escenarios de precio generados con los cuantiles previstos (aproximación por media muestral, M24, §6.3) o una versión robusta; se compara con usar la media (M24, §6.1).
  \item \emph{Control del horno} por MPC económico (M26, definición~3.10) con horizonte de 4~h y paso de 5~min sobre un modelo identificado con los datos de entrenamiento (ARX, M26, §2.1), restricciones de banda durante el tratamiento ablandadas con penalización exacta (M26, §3.3) y un filtro de seguridad de potencia y temperatura (M26, §6.5).
  \item \emph{Ajuste}: los pesos del MPC, el margen de seguridad de temperatura y la holgura de capacidad del plan se ajustan con optimización bayesiana sobre las semanas de validación del gemelo, con un máximo de 60 evaluaciones (M25, tema~4).
  \item \emph{Opcional}: un agente SAC o TD3 entrenado en el gemelo (M26, tema~4), comparado con el MPC con el mismo coste y las mismas semanas.
\end{enumerate}

\begin{table}[H]
\centering\small
\caption{Opción (a): métodos sugeridos y dónde se estudiaron.}\label{tab:metodos-a}
\begin{tabularx}{\textwidth}{@{}>{\raggedright\arraybackslash}p{3.6cm}>{\raggedright\arraybackslash}X>{\raggedright\arraybackslash}p{3.8cm}@{}}
\toprule
Componente & Método & Referencia del curso \\
\midrule
Plan semanal & PLEM con big-M mínimas, horizonte deslizante, precios sombra & M24, §2.2--§2.4; M24, §3.3 (taller) \\
Precio y pedidos & regresión cuantílica, \emph{gradient boosting}, LSTM opcional & M08, temas~2 y~6; M12, §5.4; 0.9, §1.3 \\
Decisión con incertidumbre & SAA, robusta, valor de la solución estocástica & M24, §6.1--§6.3 \\
Horno & identificación ARX, MPC económico, filtro de seguridad & M26, §2.1, §3.3, §6.5 \\
Ajuste & optimización bayesiana con EI y proceso gaussiano & M25, §3.3, §4.3, §4.6 \\
Alternativa de control & SAC o TD3 en un entorno Gymnasium & M26, §4.1--§4.4; M21, §1.3 \\
Evaluación y valor & CRN, intervalos pareados y \emph{bootstrap}, VAN, CVaR & M26, §5.1; M28, §5.1--§5.4 \\
Producción & deriva (PSI), reentrenamiento, AMFE & M28, §4.3, §5.5 \\
\bottomrule
\end{tabularx}
\end{table}

\paragraph{Tecnologías.} HiGHS (vía SciPy o \texttt{highspy}), PuLP o Pyomo para el PLEM; OR-Tools CP-SAT como alternativa para la secuenciación; CVXPY con OSQP, do-mpc o CasADi para el MPC; scikit-learn para la previsión cuantílica; BoTorch u Optuna para el ajuste; Gymnasium y Stable-Baselines3 para la alternativa de RL; MLflow, FastAPI y \texttt{pytest} para el seguimiento, el servicio y las pruebas.

\begin{table}[H]
\centering\small
\caption{Opción (a): hitos semanales.}\label{tab:hitos-a}
\begin{tabularx}{\textwidth}{@{}c>{\raggedright\arraybackslash}X@{}}
\toprule
Semana & Hito \\
\midrule
1 & Concepto de operación y requisitos; generador de precios, ambiente y pedidos; gemelo del horno y de la línea; línea base EDD con PI funcionando; partición congelada. \\
2 & Validación del gemelo (balance de energía, respuesta al escalón del horno, utilización); techo con información perfecta y mejor alternativa sencilla; documento de diseño v0. \\
3 & PLEM del plan semanal con precios reales: big-M mínimas, \emph{gap} $\le 1\,\%$, tiempos, precio sombra de la potencia contratada. \\
4 & Previsión cuantílica de precios D+2 a D+7 con validación temporal (pérdida \emph{pinball}); plan con escenarios frente a plan con la media. \\
5 & Identificación ARX del horno, MPC económico, filtro de seguridad e integración plan--consignas; documento de diseño v1. \\
6 & Optimización bayesiana de pesos y márgenes en validación; pruebas de estrés: sensor congelado, pico de precio $\times 3$ durante 6~h, disparo del horno y 20\,\% más pedidos. \\
7 & Evaluación en las 12 semanas de prueba con 5 semillas y CRN; intervalos, VAN con Monte Carlo, CVaR, AMFE y PSI de las variables de precio. \\
8 & Informe final, API del plan diario y panel (plan, temperatura, potencia, coste); demostración y reproducción completa. \\
\bottomrule
\end{tabularx}
\end{table}

\paragraph{Entregables específicos.} Además de E1 a E4: un diagrama de Gantt del plan de una semana de prueba frente a la línea base, con el precio horario superpuesto; la curva de temperatura y potencia del horno durante un día con la banda de tolerancia; la tabla de costes por partida (energía, arranques, retrasos, almacenamiento) de las cinco políticas (línea base, alternativa sencilla, optimizador, techo y, si se hace, RL).

\paragraph{Criterios de aceptación.} Los obligatorios deben cumplirse todos para aprobar el proyecto (\cref{sec:proy-calificacion}); los de excelencia suben la nota en la rúbrica.
\begin{enumerate}[label=A\arabic*.]
  \item \emph{Seguridad} (obligatorio): ninguna violación de la potencia contratada ni de la banda de $\pm5$~°C durante un tratamiento en las 12 semanas de prueba con 5 semillas; ante un disparo o un sensor congelado inyectados, el sistema pasa al respaldo en el paso siguiente.
  \item \emph{Mejora} (obligatorio): coste total al menos un 8\,\% menor que el de la línea base, con el extremo inferior del intervalo de confianza del 95\,\% pareado por semana de al menos el 4\,\%.
  \item \emph{Servicio} (obligatorio): al menos el 97\,\% de los lotes a tiempo y no más de 0{,}5 puntos por debajo de la línea base.
  \item \emph{Cómputo} (obligatorio): plan diario en menos de 120~s con \emph{gap} $\le 1\,\%$ (o el \emph{gap} alcanzado, justificado), y paso del MPC en menos de 2~s en el percentil~99.
  \item \emph{Honestidad} (obligatorio): mejora frente a la alternativa sencilla con su intervalo y fracción capturada del margen entre esta y el techo.
  \item \emph{Reproducibilidad} (obligatorio): la orden única reproduce las cifras del informe con una diferencia relativa menor del 0{,}5\,\%.
  \item \emph{Excelencia}: mejora $\ge 12\,\%$; pérdida \emph{pinball} de la previsión de D+2 al menos un 15\,\% menor que la de la previsión estacional ingenua; captura de al menos el 50\,\% del margen hasta el techo; $\Prob(\VAN<0)\le 10\,\%$.
\end{enumerate}

\begin{table}[p]
\centering\footnotesize
\caption{Rúbrica de la opción (a): planta industrial (100 puntos).}\label{tab:rubrica-a}
\begin{tabularx}{\textwidth}{@{}>{\raggedright\arraybackslash}p{2.5cm}c>{\raggedright\arraybackslash}X>{\raggedright\arraybackslash}X>{\raggedright\arraybackslash}X@{}}
\toprule
Criterio & Peso & Excelente & Suficiente & Insuficiente \\
\midrule
Documento de diseño y requisitos & 10 & Plantilla completa; requisitos con unidades y márgenes; línea base, alternativa sencilla y techo calculados en la semana~2 & Plantilla completa con alguna referencia calculada tarde & Faltan secciones, el objetivo no es escalar o no hay línea base \\
\addlinespace
Datos y gemelo digital & 10 & Generador fiel al escenario, gemelo validado con balance de energía y escalón, partición congelada con \emph{hash} & Generador correcto y gemelo validado parcialmente & Datos simplificados, fugas o prueba usada para elegir \\
\addlinespace
Plan semanal (PLEM) & 20 & Formulación completa con big-M mínimas justificadas, horizonte deslizante, \emph{gap} y tiempos, precio sombra de la potencia interpretado & Formulación correcta con alguna simplificación (por ejemplo, sin turnos) & Big-M arbitrarias, restricciones duras ausentes o plan infactible \\
\addlinespace
Control del horno & 15 & ARX validado en simulación libre, MPC económico con restricciones blandas exactas y filtro de seguridad demostrado & MPC o PI mejorado que cumple la banda & Sin control o con violaciones de la banda \\
\addlinespace
Componente de IA & 10 & Previsión cuantílica calibrada, escenarios en el plan, BO del ajuste con curva de convergencia & Previsión puntual bien validada y ajuste manual razonado & Previsión sin validación temporal \\
\addlinespace
Evaluación & 15 & CRN con 5 semillas, intervalos $t$ y \emph{bootstrap}, ablaciones y las cuatro pruebas de estrés & Intervalos correctos y alguna prueba de estrés & Solo medias o comparación no pareada \\
\addlinespace
Valor, riesgo y producción & 10 & VAN con Monte Carlo, CVaR, AMFE, PSI y API del plan & VAN determinista y AMFE & Sin análisis de valor \\
\addlinespace
Código y demostración & 10 & Orden única, pruebas automáticas, registro de experimentos, demostración con salvaguarda activada & Código reproducible con pasos manuales & No reproducible \\
\midrule
\textbf{Total} & \textbf{100} & & & \\
\bottomrule
\end{tabularx}
\end{table}

% =====================================================================
\section{Opción (b): software 3D}
\label{sec:proy-b}

\subsection*{Niveles de detalle y geometría de una escena con presupuesto de fotograma}

\paragraph{Contexto y decisiones.} Un estudio prepara una escena exterior para una consola de referencia a 60~Hz. El optimizador decide en dos niveles: (i) \emph{fuera de línea}, en la compilación de la escena, la cadena de niveles de detalle (LOD) de cada malla base (cuántos niveles, con qué fracción de triángulos y qué parámetros de simplificación) y el formato de compresión de cada textura, y (ii) \emph{en línea}, en cada fotograma, el nivel de detalle de cada instancia visible, de modo que el tiempo de fotograma quede dentro del presupuesto con el menor error visual y sin parpadeos de cambio de nivel (\emph{popping}).

\paragraph{Datos sintéticos que hay que generar.}
\begin{enumerate}[label=(\roman*)]
  \item \emph{Biblioteca de 40 mallas base} generadas de forma procedural con semilla: rocas (icoesferas subdivididas cinco veces con desplazamiento por ruido fractal), árboles (tronco cilíndrico y copas de icoesferas deformadas), edificios (cajas con ventanas extruidas y cornisas) y personajes (superficie implícita de esferas y cápsulas extraída con \emph{marching cubes}, M27, §1.3). Entre 5000 y 80\,000 triángulos por malla, cerradas y orientadas de forma coherente.
  \item \emph{Escena}: terreno de $1\times 1$~km con 2000 instancias (60\,\% vegetación, 25\,\% rocas, 12\,\% edificios, 3\,\% personajes), colocadas con un proceso de Poisson agrupado; importancia visual $w$ de 1; 0{,}5; 1{,}5 y~3 por clase.
  \item \emph{Cámaras}: seis trayectorias de 60~s a 60~fotogramas por segundo (3600 fotogramas cada una), a pie, en vehículo y en vuelo bajo: dos para entrenamiento, dos para validación y dos para prueba. Campo de visión vertical de $60^{\circ}$ e imagen de $1920\times 1080$.
  \item \emph{Modelo de coste} (si no se mide en hardware real): $t_{\mathrm{GPU}}=5{,}0~\text{ms}+\sum_{i\ \text{visible}}\bigl(\alpha\,\mathrm{tri}_i+\beta\,\mathrm{px}_i+\gamma\bigr)+\epsilon$, con $\alpha=2{,}5$~ns por triángulo, $\beta=1{,}2$~ns por píxel sombreado (con un \emph{overdraw} de $1{,}6$), $\gamma=4~\mu$s por llamada de dibujo, ruido multiplicativo log-normal de desviación $0{,}05$ y, con probabilidad $0{,}2\,\%$ por fotograma, un tirón de 8~ms; $t_{\mathrm{CPU}}=2~\text{ms}+6~\mu\text{s}$ por instancia visible. Tiempo de fotograma $\max(t_{\mathrm{CPU}},t_{\mathrm{GPU}})$ (M27, definición~5.1). Medir el coste real con un renderizador (por ejemplo, \texttt{moderngl} o \texttt{pyrender}) es opcional y puntúa en la rúbrica.
  \item \emph{Calidad}: error en pantalla de cada instancia a partir del error geométrico de su nivel (distancia de Hausdorff a LOD0) con la proposición~5.5 de M27, ponderado por la importancia; métrica perceptual opcional (por ejemplo, FLIP o SSIM) sobre imágenes renderizadas de un subconjunto de fotogramas.
  \item \emph{Memoria}: tres texturas de $4096\times 4096$ por malla base (color, normales y material). Color y material en BC7 (1~byte por téxel) o BC1 (0{,}5~bytes por téxel), normales en BC5 (1~byte por téxel), más un tercio para los mipmaps (M27, §5.5); se puede además descartar el primer nivel de mipmap de una textura, lo que divide su memoria por cuatro. Todo en BC7 y BC5 a resolución completa ocupa unos 2{,}7~GB, frente a un presupuesto de memoria de vídeo de 1{,}5~GB para texturas y geometría (geometría a unos 36~bytes por vértice, M27, §1.2): la elección de formatos y resoluciones es, pues, una segunda mochila de elección múltiple.
\end{enumerate}

\paragraph{Referencias de comparación.} \emph{Línea base}: umbrales fijos de distancia de 10, 30 y 80~m para pasar a LOD1, LOD2 y LOD3, con LOD generados con una fracción fija de triángulos (50, 25 y 10\,\%). \emph{Mejor alternativa sencilla}: umbrales por error proyectado de 1~píxel (M27, §5.3), sin presupuesto. \emph{Techo}: el PLEM por ventanas de 60 fotogramas con la trayectoria conocida y los costes exactos.

\paragraph{Formulación que se pide.}
\begin{enumerate}[label=F\arabic*.]
  \item \emph{Cadena de LOD} por malla base: simplificación por error cuádrico (M27, §5.2) con fracciones de triángulos y opciones (preservar bordes, peso de las normales) elegidas para minimizar el error en pantalla esperado en las trayectorias de entrenamiento con un límite de memoria; los parámetros se ajustan con optimización bayesiana (M25, tema~4) o con halving sucesivo si la evaluación es cara (M28, §3.3). Cada nivel pasa la validación geométrica (M27, §6.2).
  \item \emph{Selección por fotograma} como mochila de elección múltiple (M27, definición~5.6): $\min\sum_iw_i\,e_{i\ell}(d_i)+\kappa\sum_i\mathbb{1}[\ell_{i,t}\ne\ell_{i,t-1}]$ con $\sum_it_{i\ell}\le B_t$, cambios de un solo nivel por fotograma y la regla de dirección de arte de que los personajes a menos de 15~m van en LOD0. Se resuelve por relajación lagrangiana (M27, §5.4), y el multiplicador $\lambda$ se actualiza con un controlador PI sobre el tiempo de fotograma medido (M26, §1.4) para cumplir el percentil~99.
  \item \emph{Modelo de coste aprendido}: regresión (\emph{gradient boosting} o perceptrón multicapa, M08 y M09) del tiempo en función de triángulos, píxeles, llamadas de dibujo y material, validada por trayectorias y usada dentro de F2.
  \item \emph{Techo}: PLEM sobre ventanas de 60 fotogramas con el coste exacto, para acotar lo alcanzable (M24, §2.3).
  \item \emph{Opcional}: búsqueda de la función de prioridad de la selección con un LLM (FunSearch u OpenEvolve, M28, tema~1), con el simulador de fotogramas como evaluador y las trayectorias de validación separadas de las de búsqueda.
\end{enumerate}

\begin{table}[H]
\centering\small
\caption{Opción (b): métodos sugeridos y dónde se estudiaron.}\label{tab:metodos-b}
\begin{tabularx}{\textwidth}{@{}>{\raggedright\arraybackslash}p{3.6cm}>{\raggedright\arraybackslash}X>{\raggedright\arraybackslash}p{3.8cm}@{}}
\toprule
Componente & Método & Referencia del curso \\
\midrule
Mallas y validación & mallas, característica de Euler, SDF, comprobaciones topológicas & M27, §1.2, §1.3 y §6.2 \\
Simplificación & error cuádrico, error proyectado & M27, §5.2 y §5.3 \\
Selección con presupuesto & mochila de elección múltiple, relajación lagrangiana, PLEM & M27, §5.4; M24, §2.3 y §4.2 \\
Control del presupuesto & PI sobre el multiplicador, antisaturación & M26, §1.4 \\
Ajuste de la cadena & optimización bayesiana, halving sucesivo, hipercubo latino & M25, §2.6 y tema~4; M28, §3.3 \\
Modelo de coste & regresión, \emph{gradient boosting}, perceptrón multicapa & M08, temas~2 y~6; M09, tema~1 \\
Métrica perceptual & renderizado diferenciable y comparación de imágenes & M27, §4.1 y §4.2 \\
Búsqueda de heurísticas & FunSearch, AlphaEvolve, sesgo del ganador & M28, §1.2--§1.4 y §2.2 \\
Evaluación y valor & \emph{bootstrap} por bloques, VAN, monitorización & M28, §5.1--§5.2, §4.3 \\
\bottomrule
\end{tabularx}
\end{table}

\paragraph{Tecnologías.} Blender con \texttt{bpy} (modificador \emph{Decimate} en modo de colapso) o \texttt{trimesh} y \texttt{pymeshlab} para generar y simplificar; meshoptimizer para la simplificación y la reordenación de índices; Open3D para la distancia de Hausdorff; \texttt{nvdiffrast}, PyTorch3D, \texttt{pyrender} o \texttt{moderngl} para renderizar; OR-Tools o HiGHS para el PLEM del techo; scikit-learn y PyTorch para el modelo de coste; Optuna o BoTorch para el ajuste; OpenEvolve para la búsqueda opcional; MLflow y una interfaz de línea de órdenes integrada en la compilación de la escena.

\begin{table}[H]
\centering\small
\caption{Opción (b): hitos semanales.}\label{tab:hitos-b}
\begin{tabularx}{\textwidth}{@{}c>{\raggedright\arraybackslash}X@{}}
\toprule
Semana & Hito \\
\midrule
1 & Concepto de operación y requisitos (percentil 99, píxeles, memoria); generador de mallas y escena; trayectorias; simulador de fotogramas; línea base por distancias; partición congelada. \\
2 & Validación geométrica de la biblioteca; modelo de coste calibrado; techo por ventanas y alternativa por error proyectado; documento de diseño v0. \\
3 & Cadenas de LOD por error cuádrico con validación; curvas error--triángulos por malla. \\
4 & Modelo de coste aprendido con validación por trayectorias; selección por fotograma con relajación lagrangiana y $\lambda$ fijo. \\
5 & Control PI de $\lambda$, histéresis y regla de dirección de arte; formato de texturas con el presupuesto de memoria; documento de diseño v1. \\
6 & Optimización bayesiana de los parámetros de simplificación y del controlador en validación; pruebas de estrés: giro rápido de cámara, multitud de personajes, GPU un 30\,\% más lenta y tirones de compilación de sombreadores. \\
7 & Evaluación en las dos trayectorias de prueba con 5 semillas del ruido del coste; intervalos \emph{bootstrap} por bloques de 1~s; comprobación perceptual opcional. \\
8 & Informe, herramienta integrada en la compilación, vídeo comparativo lado a lado y demostración. \\
\bottomrule
\end{tabularx}
\end{table}

\paragraph{Entregables específicos.} Además de E1 a E4: un vídeo de 30~s con la línea base y el optimizador lado a lado y el tiempo de fotograma superpuesto; la gráfica del tiempo de fotograma y de $\lambda$ a lo largo de una trayectoria de prueba; el informe de validación geométrica de todos los LOD generados.

\paragraph{Criterios de aceptación.}
\begin{enumerate}[label=B\arabic*.]
  \item \emph{Rendimiento} (obligatorio): percentil~99 del tiempo de fotograma $\le 16{,}7$~ms en las dos trayectorias de prueba y ninguna ventana de 1~s con media superior a $16{,}7$~ms.
  \item \emph{Calidad} (obligatorio): error en pantalla ponderado medio al menos un 20\,\% menor que el de la línea base con un percentil~99 igual o menor, con un intervalo \emph{bootstrap} del 95\,\% por bloques de 1~s cuyo extremo inferior sea de al menos el 10\,\%.
  \item \emph{Estabilidad} (obligatorio): como máximo un cambio de nivel por instancia visible cada 2~s de media, y ninguna instancia con más de 3 cambios en 1~s.
  \item \emph{Cómputo y memoria} (obligatorio): selección de las 2000 instancias en menos de 0{,}5~ms en el percentil~99 en CPU; memoria de vídeo dentro del presupuesto en todos los fotogramas.
  \item \emph{Validez geométrica} (obligatorio): el 100\,\% de los LOD generados son variedades cerradas con orientación coherente, sin caras degeneradas, con la característica de Euler de su malla base y con un error de Hausdorff no mayor que el declarado.
  \item \emph{Honestidad y reproducibilidad} (obligatorio): mejora frente a la alternativa por error proyectado con su intervalo, fracción capturada del margen hasta el techo, y reproducción con una sola orden.
  \item \emph{Excelencia}: error del modelo de coste (MAPE) $\le 10\,\%$ en prueba; correlación de Spearman $\ge 0{,}7$ entre el error en pantalla y la métrica perceptual en 200 fotogramas; captura de al menos el 60\,\% del margen hasta el techo; coste medido en hardware real.
\end{enumerate}

\begin{table}[p]
\centering\footnotesize
\caption{Rúbrica de la opción (b): software 3D (100 puntos).}\label{tab:rubrica-b}
\begin{tabularx}{\textwidth}{@{}>{\raggedright\arraybackslash}p{2.5cm}c>{\raggedright\arraybackslash}X>{\raggedright\arraybackslash}X>{\raggedright\arraybackslash}X@{}}
\toprule
Criterio & Peso & Excelente & Suficiente & Insuficiente \\
\midrule
Documento de diseño y requisitos & 10 & Requisitos en percentil~99, píxeles y memoria con márgenes; las tres referencias calculadas en la semana~2 & Plantilla completa con alguna referencia tardía & Requisitos sobre la media o sin línea base \\
\addlinespace
Datos y modelo de coste & 10 & Biblioteca y escena fieles al escenario; modelo de coste validado por trayectorias; partición congelada & Generador correcto, coste sin validar por trayectorias & Escena trivial o trayectorias de prueba usadas para ajustar \\
\addlinespace
Cadena de LOD & 15 & Error cuádrico con parámetros ajustados por BO o halving, curvas error--triángulos, validación geométrica completa & Simplificación correcta con parámetros fijos razonados & LOD inválidos o sin validar \\
\addlinespace
Selección con presupuesto & 20 & Mochila con relajación lagrangiana, cota dual, control PI de $\lambda$, histéresis y regla de arte; techo por PLEM & Lagrangiana con $\lambda$ fijo ajustado a mano & Selección voraz sin presupuesto o que lo incumple \\
\addlinespace
Componente de IA & 10 & Modelo de coste aprendido integrado y validado; búsqueda con LLM o métrica perceptual & Modelo de coste aprendido sin integrar & Ninguno o sin validar \\
\addlinespace
Evaluación & 15 & \emph{Bootstrap} por bloques, 5 semillas, ablaciones (sin histéresis, sin control de $\lambda$, sin modelo aprendido), cuatro pruebas de estrés & Intervalos correctos y alguna prueba de estrés & Medias sin intervalos o en trayectorias de entrenamiento \\
\addlinespace
Integración y valor & 10 & Herramienta en la compilación, telemetría del tiempo de fotograma, horas de artista ahorradas y VAN & Script reproducible con informe de memoria & Proceso manual \\
\addlinespace
Código y demostración & 10 & Orden única, pruebas automáticas, vídeo comparativo y demostración con caso de estrés & Código reproducible con pasos manuales & No reproducible \\
\midrule
\textbf{Total} & \textbf{100} & & & \\
\bottomrule
\end{tabularx}
\end{table}

% =====================================================================
\section{Opción (c): operaciones}
\label{sec:proy-c}

\subsection*{Rutas y mantenimiento predictivo de una flota de reparto}

\paragraph{Contexto y decisiones.} Una empresa de reparto urbano opera 30 furgonetas desde un depósito; cada una recorre unos 100~km por día laborable. El optimizador toma tres decisiones acopladas: (i) cada domingo, y con revisión diaria, el \emph{plan de mantenimiento}: qué vehículo pasa por el taller qué día, con una capacidad de dos vehículos diarios de lunes a viernes; (ii) cada noche, las \emph{rutas} del día siguiente con los vehículos disponibles, respetando capacidades, ventanas horarias y jornada, y (iii) la \emph{asignación según el riesgo}: los vehículos con mayor probabilidad de avería reciben rutas más cortas y cercanas al depósito, porque una avería lejos cuesta más.

\paragraph{Datos sintéticos que hay que generar.} Dos años ficticios de operación diaria.
\begin{enumerate}[label=(\roman*)]
  \item \emph{Ciudad}: cuadrado de $30\times 30$~km con el depósito en $(5,15)$~km; 4000 direcciones en cinco agrupaciones gaussianas. Tiempo de viaje: distancia de Manhattan por $1{,}3$ dividida por una velocidad de 28~km/h, que baja a 18~km/h de 8:00 a 9:30 y de 17:30 a 19:00, con ruido multiplicativo log-normal de desviación $0{,}15$.
  \item \emph{Pedidos}: cada día laborable, una Poisson de media 650 (factor $1{,}2$ los lunes y en diciembre), de 1 a 4 bultos, con servicio de $4+0{,}5$ minutos por bulto; ventanas de mañana (8--12~h, 40\,\%), de tarde (14--18~h, 35\,\%) o libres (8--18~h, 25\,\%).
  \item \emph{Flota y costes}: capacidad de 70 bultos (con la demanda media hacen falta unas 23 furgonetas por capacidad, y los lunes, unas 28); jornada de 8 a 18~h ampliable hasta las 20~h con horas extra a 35~€/h; conductor a 22~€/h; 0{,}25~€/km; entrega fuera de ventana, 15~€; subcontratar un vehículo un día, 250~€.
  \item \emph{Degradación}: un componente crítico por vehículo con riesgos proporcionales, $\lambda(t\mid z)=\lambda_0(t)\exp(0{,}6\,z_{\mathrm{carga}}+0{,}4\,z_{\mathrm{paradas}})$, con $\lambda_0$ de Weibull de forma $\beta=2{,}5$ y escala $\theta=15\,000$~km (M28, definición~4.11), y covariables tipificadas que dependen de las rutas asignadas. El fallo se produce cuando el riesgo acumulado supera una exponencial de media~1 sorteada tras cada renovación. Telemetría diaria: índice de vibración $v=1+2D^{2}+\N(0,\,0{,}15^{2})$, con $D\in[0,1]$ el cociente entre el riesgo acumulado y el umbral sorteado (la degradación), temperatura del componente y códigos de error de Poisson con tasa creciente en $D$.
  \item \emph{Mantenimiento}: preventivo, 350~€ y un día sin vehículo; avería en ruta, 1800~€ (grúa y reparación), los pedidos pendientes se reasignan con penalización y el vehículo queda tres días fuera. El mantenimiento renueva el componente («como nuevo»).
  \item \emph{Histórico y registro}: la política histórica revisa cada vehículo a los 10\,000~km y, además, cada semana inspecciona al azar un vehículo con probabilidades conocidas y registradas (para la evaluación contrafactual, M28, §4.4). Las vidas interrumpidas por el preventivo quedan censuradas (M28, observación~4.12). En el mes~21 se sustituye el sensor de vibración y su lectura pasa a tener un sesgo de $+0{,}3$.
  \item \emph{Partición}: meses 1 a~15 para entrenamiento, 16 a~19 para validación y 20 a~24 para prueba.
\end{enumerate}

\paragraph{Referencias de comparación.} \emph{Línea base}: revisión cada 10\,000~km y rutas con la construcción por el arco más barato de OR-Tools sin búsqueda local. \emph{Mejor alternativa sencilla}: la edad óptima de sustitución $T^{\ast}$ calculada con una Weibull ajustada con censura (M28, proposición~4.13) y rutas con búsqueda local guiada durante 30~s. \emph{Techo}: un oráculo que conoce el instante de cada fallo y revisa el día anterior, con los tiempos de viaje reales.

\paragraph{Formulación que se pide.}
\begin{enumerate}[label=F\arabic*.]
  \item \emph{Modelo de riesgo}: probabilidad de fallo en los próximos 7 días por vehículo, con un modelo de supervivencia de Weibull con riesgos proporcionales estimado por máxima verosimilitud con censura y con un clasificador (\emph{gradient boosting} o regresión logística, M08, tema~3) con validación temporal; se comparan discriminación (AUC), calibración (error de calibración esperado con 10 intervalos) y, sobre todo, el coste de las decisiones que inducen.
  \item \emph{Plan de mantenimiento semanal} como PLEM de asignación: $m_{vd}\in\{0,1\}$ (el vehículo $v$ va al taller el día $d$, o no va, $d=\varnothing$), $\sum_vm_{vd}\le 2$, $\sum_dm_{vd}\le 1$; coste esperado de cada opción $c_{vd}=c_p+c_f\sum_{d'<d}\hat r_{vd'}+\text{coste de capacidad}$, donde $\hat r_{vd'}$ es el riesgo diario previsto. El coste de capacidad penaliza los días en que los vehículos disponibles no cubren la demanda prevista (subcontratación).
  \item \emph{Rutas diarias} como VRPTW (M24, §3.2) resuelto con OR-Tools (dimensiones de tiempo y de carga), objetivo de kilómetros, horas de conductor, horas extra y entregas fuera de ventana, y un coste por kilómetro que crece con el riesgo del vehículo, $0{,}25+\kappa\,\hat r_v$~€/km, con $\kappa$ elegido en validación.
  \item \emph{Evaluación contrafactual}: antes de simular, estimar con ponderación inversa de la propensión el valor de la nueva política de inspección a partir de las inspecciones aleatorias del histórico (M28, proposición~4.9), y compararlo con el gemelo.
  \item \emph{Producción}: PSI semanal de la telemetría con umbrales calibrados en un periodo estable, periodo de reentrenamiento (M28, proposición~4.8) y AMFE del sistema.
\end{enumerate}

\begin{table}[H]
\centering\small
\caption{Opción (c): métodos sugeridos y dónde se estudiaron.}\label{tab:metodos-c}
\begin{tabularx}{\textwidth}{@{}>{\raggedright\arraybackslash}p{3.6cm}>{\raggedright\arraybackslash}X>{\raggedright\arraybackslash}p{3.8cm}@{}}
\toprule
Componente & Método & Referencia del curso \\
\midrule
Rutas & VRPTW, big-M de las ventanas, búsqueda local, 2-opt & M24, §3.1 y §3.2 \\
Plan de mantenimiento & PLEM de asignación con capacidad, precios sombra del taller & M24, §2.2 y §2.3 \\
Fiabilidad & Weibull, censura, política por edad & M28, §4.5; 0.9, §4.2 \\
Modelo de riesgo & regresión logística, \emph{gradient boosting}, calibración & M08, tema~3; M05, §4.2 \\
Decisión con incertidumbre & coste esperado, programación estocástica, CVaR & M24, §6.3; M28, §5.4 \\
Evaluación contrafactual & ponderación inversa de la propensión & M28, §4.4 \\
Ajuste de $\kappa$ y del modelo & optimización bayesiana, búsqueda aleatoria & M25, tema~4 y §6.2 \\
Producción y valor & PSI, reentrenamiento, intervalos, VAN, AMFE & M28, §4.3, §5.1--§5.5 \\
\bottomrule
\end{tabularx}
\end{table}

\paragraph{Tecnologías.} OR-Tools (rutas y CP-SAT) o PyVRP para el VRPTW; HiGHS o PuLP para el plan de mantenimiento; \texttt{lifelines} o \texttt{scikit-survival} para la supervivencia con censura, o una verosimilitud propia con SciPy; scikit-learn para los clasificadores y la calibración; Optuna para el ajuste; Evidently o un PSI propio para la deriva; MLflow, FastAPI y \texttt{pytest}; un mapa con \texttt{folium} o \texttt{matplotlib} para las rutas.

\begin{table}[H]
\centering\small
\caption{Opción (c): hitos semanales.}\label{tab:hitos-c}
\begin{tabularx}{\textwidth}{@{}c>{\raggedright\arraybackslash}X@{}}
\toprule
Semana & Hito \\
\midrule
1 & Concepto de operación y requisitos; generador de ciudad, pedidos, tiempos y degradación; gemelo de un día de operación; línea base; partición congelada. \\
2 & Validación del gemelo (kilómetros, puntualidad, tasa de averías de la política histórica); techo con oráculo y política $T^{\ast}$; documento de diseño v0. \\
3 & VRPTW diario con OR-Tools: tiempos, puntualidad, comparación de estrategias de búsqueda; certificación de un día pequeño con un PLEM o con programación dinámica. \\
4 & Modelos de riesgo con censura y clasificador; validación temporal, AUC y calibración; estimación IPS de la política de inspección. \\
5 & PLEM del plan de mantenimiento, acoplamiento con las rutas (disponibilidad y coste por riesgo); documento de diseño v1. \\
6 & Ajuste de $\kappa$ y de los umbrales en validación; pruebas de estrés: pico de demanda del 30\,\%, taller con un solo puesto una semana, sensor congelado y lluvia (velocidades un 20\,\% menores). \\
7 & Evaluación en los 5 meses de prueba con 5 semillas y CRN; intervalos por semana, VAN, CVaR, AMFE; detección del cambio de sensor del mes~21. \\
8 & Informe, API de planificación nocturna, panel de riesgo y rutas, demostración. \\
\bottomrule
\end{tabularx}
\end{table}

\paragraph{Entregables específicos.} Además de E1 a E4: el mapa de las rutas de un día de prueba con los vehículos coloreados por riesgo; la curva de fiabilidad estimada con censura frente a la verdadera; el calendario de taller de una semana; la serie del PSI alrededor del cambio de sensor.

\paragraph{Criterios de aceptación.}
\begin{enumerate}[label=C\arabic*.]
  \item \emph{Restricciones duras} (obligatorio): capacidad del taller, capacidad de los vehículos y jornada máxima (20:00) siempre respetadas; todo pedido se asigna o se subcontrata de forma explícita.
  \item \emph{Mejora} (obligatorio): coste total (kilómetros, conductores, horas extra, retrasos, mantenimiento, averías y subcontratación) al menos un 8\,\% menor que el de la línea base en los 5 meses de prueba, con el extremo inferior del intervalo del 95\,\% pareado por semana de al menos el 4\,\%.
  \item \emph{Servicio y averías} (obligatorio): al menos el 96\,\% de entregas en ventana y al menos un 30\,\% menos de averías en ruta que la línea base.
  \item \emph{Modelo de riesgo} (obligatorio): AUC $\ge 0{,}80$ y error de calibración esperado $\le 0{,}03$ en la prueba temporal; el estimador de Weibull con censura recupera $\beta$ y $\theta$ con un error menor del 10\,\% en datos simulados con parámetros conocidos.
  \item \emph{Cómputo y monitorización} (obligatorio): rutas del día en menos de 120~s y plan de mantenimiento en menos de 10~s; alarma de deriva en los 7 días siguientes al cambio de sensor, con no más de una falsa alarma por trimestre en el periodo estable.
  \item \emph{Honestidad y reproducibilidad} (obligatorio): mejora frente a la alternativa $T^{\ast}$ con su intervalo, aunque contenga el cero; fracción capturada del margen hasta el techo; reproducción con una sola orden.
  \item \emph{Excelencia}: captura de al menos el 40\,\% del margen entre la alternativa sencilla y el techo; estimación IPS dentro del intervalo del gemelo; $\Prob(\VAN<0)\le 10\,\%$.
\end{enumerate}

\begin{table}[p]
\centering\footnotesize
\caption{Rúbrica de la opción (c): operaciones (100 puntos).}\label{tab:rubrica-c}
\begin{tabularx}{\textwidth}{@{}>{\raggedright\arraybackslash}p{2.5cm}c>{\raggedright\arraybackslash}X>{\raggedright\arraybackslash}X>{\raggedright\arraybackslash}X@{}}
\toprule
Criterio & Peso & Excelente & Suficiente & Insuficiente \\
\midrule
Documento de diseño y requisitos & 10 & Decisiones acopladas bien delimitadas; requisitos con márgenes; las tres referencias, incluida $T^{\ast}$, en la semana~2 & Plantilla completa con alguna referencia tardía & Sin línea base o con la práctica actual como única referencia \\
\addlinespace
Datos y gemelo & 10 & Generador fiel, censura y registro de propensiones correctos, gemelo validado, partición congelada & Generador correcto y gemelo parcialmente validado & Sin censura, fugas o prueba usada para ajustar \\
\addlinespace
Rutas (VRPTW) & 15 & OR-Tools con ventanas y capacidad, comparación de estrategias, certificación en un caso pequeño, coste por riesgo & VRPTW correcto sin acoplamiento con el riesgo & Rutas infactibles o sin ventanas \\
\addlinespace
Plan de mantenimiento & 15 & PLEM de asignación con capacidad y coste de capacidad, precios sombra del taller, horizonte deslizante & Asignación correcta por reglas razonadas & Sin plan o incumple la capacidad \\
\addlinespace
Componente de IA & 15 & Supervivencia con censura y clasificador calibrado comparados por el coste de decisión; IPS & Un modelo bien validado y calibrado & Sin validación temporal o sin censura \\
\addlinespace
Evaluación & 15 & CRN con 5 semillas, intervalos por semana, ablaciones (sin riesgo en las rutas, sin plan) y cuatro pruebas de estrés & Intervalos correctos y alguna prueba de estrés & Solo medias o comparación no pareada \\
\addlinespace
Valor, riesgo y producción & 10 & VAN con Monte Carlo, CVaR, AMFE, PSI con detección del cambio de sensor, API & VAN determinista y PSI & Sin análisis de valor ni monitorización \\
\addlinespace
Código y demostración & 10 & Orden única, pruebas, registro, demostración con avería inyectada & Reproducible con pasos manuales & No reproducible \\
\midrule
\textbf{Total} & \textbf{100} & & & \\
\bottomrule
\end{tabularx}
\end{table}

% =====================================================================
\section{Calificación del proyecto}
\label{sec:proy-calificacion}

\begin{enumerate}[label=\arabic*.]
  \item El proyecto se puntúa sobre 100 con la rúbrica de la opción elegida (\cref{tab:rubrica-a,tab:rubrica-b,tab:rubrica-c}). Dentro de cada criterio, «Excelente» corresponde al 90--100\,\% de su peso, «Suficiente» al 60--75\,\% e «Insuficiente» al 0--40\,\%.
  \item Se aprueba con 70 puntos o más y con todos los criterios de aceptación obligatorios cumplidos. Si un criterio obligatorio no se cumple, la nota no puede superar 65 puntos, salvo que el informe demuestre con datos que el criterio es inalcanzable con el escenario nominal (por ejemplo, porque el techo está por debajo del umbral) y proponga y justifique uno revisado antes de la semana~6; en ese caso decide el profesor.
  \item Cada criterio de excelencia cumplido y bien documentado puede compensar hasta 2 puntos perdidos en otros criterios, sin superar 100.
  \item Se penaliza con 10 puntos cada semana de retraso del documento de diseño v0 o v1, porque esas revisiones existen para corregir el rumbo a tiempo.
  \item La nota del hito combina examen y proyecto como se indica en las instrucciones: $0{,}4\cdot\text{examen}/10+0{,}6\cdot\text{proyecto}/10$, con al menos un 7 en cada parte.
\end{enumerate}


\appendix
\renewcommand{\thechapter}{\Roman{chapter}}
% =====================================================================
\chapter{Soluciones}
\label{ap:soluciones}
% =====================================================================

Los resultados finales se dan con tres cifras significativas y los intermedios con cuatro, para que el redondeo no se acumule. Las remisiones «M24, §2.3» indican el módulo y la sección de sus apuntes, comprobadas en sus índices. Al final de cada solución figuran, entre corchetes, los criterios de corrección.

% ---------------------------------------------------------------------
\section{Parte A: clave de respuestas}

\begin{table}[h]
\centering\footnotesize
\caption{Clave de la Parte~A.}\label{tab:clave}
\begin{tabularx}{\textwidth}{@{}c c >{\raggedright\arraybackslash}X l@{}}
\toprule
N.º & Resp. & Justificación & Repasar \\
\midrule
1 & (c) & La menor cota válida es el máximo factible de $x_t$: $\min(120,\,90)=90$~t (proposición~2.11 de M24). & M24, §2.3 \\
2 & (a) & Holgura complementaria: $y_i\,(b_i-(Ax)_i)=0$ y la holgura es 12~h, luego $y_i^{\ast}=0$. & M24, §2.2 \\
3 & (d) & Fractil $30/40=0{,}75$; $q^{\ast}=200+40\cdot 0{,}6745=227{,}0$~hL. 234 usa el fractil $0{,}8$. & M24, §6.2 \\
4 & (b) & Pesos positivos dan puntos Pareto-óptimos; los no apoyados no se alcanzan (proposición~5.5). & M24, §5.4 \\
5 & (c) & $z=0$: $\EI=0\cdot\Phi(0)+2\,\varphi(0)=2\cdot 0{,}3989=0{,}798$; la EI es positiva aunque $\mu=\fmejor$. & M25, §4.3 \\
6 & (a) & $\sigma^{2}=1-k_\ast^{2}=1-(e^{-1/2})^{2}=1-e^{-1}=0{,}632$. & M25, §3.3 \\
7 & (d) & $n\ge\ln 0{,}05/\ln 0{,}95=58{,}4$, luego $n=59$. & M25, §6.2 \\
8 & (b) & Con la varianza predictiva, $\sigma$ no baja de $\sigma_\varepsilon$ en los puntos ensayados (observación~4.8). & M25, §4.5 \\
9 & (a) & Autovalor de Euler $1-5\Delta t$; $\abs{1-5\Delta t}<1\iff\Delta t<2/5=0{,}4$~h. & M26, §1.3 \\
10 & (c) & $\min_u\{u^{2}+(2+u)^{2}\}$: $2u+2(2+u)=0$, $u=-1$; $K_0=P_1/(R+P_1)=1/2$. & M26, §3.1 \\
11 & (d) & Penalización exacta: basta $\rho\ge\max_i\lambda_i^{\ast}=3$ (proposición~3.9). & M26, §3.3 \\
12 & (b) & Truncar no es terminar: $d=\texttt{terminated}=0$ y $y=r+\gamma Q_{\phi'}(s',\tilde a')$. & M26, §4.4 \\
13 & (c) & $E_e/E_0=x_e^{p}=0{,}6^{3}=0{,}216$. & M27, §2.2 \\
14 & (a) & La flexibilidad es autoadjunta: $\lambda=U$ y $\partial c/\partial x_e=-p\,x_e^{p-1}E_0\,u_e\T k_0u_e$. & M27, §2.3 \\
15 & (d) & $d^{\ast}=\varepsilon H/(2\tau\tan 30^{\circ})=0{,}02\cdot 1080/1{,}1547=18{,}7$~m. & M27, §5.3 \\
16 & (a) & Corolario~1.4: $\abs{F}=2\abs{V}-4=1996$ (y $\abs{E}=2994$). & M27, §1.2 \\
17 & (b) & $\sigma\sqrt{2\ln K}=0{,}04\sqrt{2\ln 50}=0{,}04\cdot 2{,}797=0{,}112$. & M28, §2.2 \\
18 & (d) & $0{,}1\ln 1{,}2+(-0{,}1)\ln 0{,}8=0{,}01823+0{,}02231=0{,}0405$. & M28, §4.3 \\
19 & (c) & $(K+1)/\eta^{K}=4/27=0{,}148$. & M28, §3.3 \\
20 & (b) & $a_{4;\,0{,}1}=3{,}1699$; $\VAN=20\,000\cdot 3{,}1699-50\,000=13\,397$~€. & M28, §5.1 \\
\bottomrule
\end{tabularx}
\end{table}

% ---------------------------------------------------------------------
\section{Parte B: soluciones de los problemas}

% ------------------------------------------------------------- B.1
\begin{solucion}{eje:milp}
\datos $D=100$~t; $f=(400,150,250)$~€; $c=(12,16,14)$~€/t; $L=(30,20,25)$~t; $U=(80,50,60)$~t.
\planteamiento PLEM con costes fijos y lotes mínimos (M24, §2.3); la menor $M$ válida es el máximo factible de $x_i$ (proposición~2.11 de M24); la relajación lineal acota el óptimo (proposición~2.10) y la disyunción se escribe con una binaria (M24, §2.3, apartado~(d) de la lista de modelizaciones lógicas).
\pasos

\noindent\textbf{(a)}
\begin{enumerate}[label=\arabic*.]
  \item Modelo:
  \[
  \begin{aligned}
  \min\ &\sum_{i=1}^{3}\bigl(f_iy_i+c_ix_i\bigr)\\
  \text{s.a.}\ &x_1+x_2+x_3=100,\qquad L_iy_i\le x_i\le M_iy_i,\qquad x_i\ge 0,\quad y_i\in\{0,1\}\quad(i=1,2,3).
  \end{aligned}
  \]
  \item La igualdad cubre exactamente la carga del turno. $x_i\le M_iy_i$ es el enlace: si el horno está apagado ($y_i=0$) no trata nada y, si trata algo, se paga $f_i$. $L_iy_i\le x_i$ es el lote mínimo: o $x_i=0$ o $x_i\ge L_i$.
\end{enumerate}

\noindent\textbf{(b)}
\begin{enumerate}[label=\arabic*.]
  \item En el conjunto factible sin el enlace, $x_i\le U_i$ (capacidad) y $x_i\le D=100$ (la igualdad con $x\ge 0$). Luego $\sup x_i=\min(U_i,D)$: $M=(80,\,50,\,60)$~t.
  \item Con $y_i=1$ la restricción es $x_i\le U_i$, que no corta nada factible; con $y_i=0$ fuerza $x_i=0$, como se desea. Una $M_i$ menor cortaría planes factibles; una mayor es válida pero debilita la relajación (proposición~2.11).
\end{enumerate}

\noindent\textbf{(c)}
\begin{enumerate}[label=\arabic*.]
  \item Un horno solo no basta: $U_i<100$ para todo $i$. Los patrones $(0,0,0)$, $(1,0,0)$, $(0,1,0)$ y $(0,0,1)$ son infactibles.
  \item $\{1,2\}$: $x_1+x_2=100$, $x_1\le 80$, $x_2\ge 20$. El horno~1 es más barato por tonelada, así que $x_1=80$, $x_2=20$: $550+960+320=1830$~€.
  \item $\{1,3\}$: $x_3\ge 25$ obliga a $x_1\le 75$: $x_1=75$, $x_3=25$: $650+900+350=1900$~€.
  \item $\{2,3\}$: el horno~3 es más barato: $x_3=60$, $x_2=40$: $400+840+640=1880$~€.
  \item $\{1,2,3\}$: mínimos $x_2=20$, $x_3=25$ y $x_1=55\le 80$: $800+660+320+350=2130$~€.
\end{enumerate}

\noindent\textbf{(d)}
\begin{enumerate}[label=\arabic*.]
  \item En la relajación conviene $y_i=x_i/M_i$ (el menor valor permitido; cumple $L_iy_i\le x_i$ porque $L_i\le M_i$). El problema se reduce a $\min\sum_i(c_i+f_i/M_i)x_i$ con $\sum_ix_i=100$, $0\le x_i\le M_i$.
  \item Costes efectivos: $12+400/80=17$, $16+150/50=19$ y $14+250/60=18{,}17$~€/t. Se llena el horno~1 (80~t) y después el~3 (20~t): $z_{\mathrm{LP}}=80\cdot 17+20\cdot 18{,}17=1723{,}3$~€, con $y=(1;\ 0;\ 1/3)$.
  \item \emph{Gap}: $(1830-1723{,}3)/1830=5{,}83\,\%$.
  \item Con $M=10^{4}$: costes efectivos $12{,}04$, $16{,}015$ y $14{,}025$~€/t; $z_{\mathrm{LP}}=80\cdot 12{,}04+20\cdot 14{,}025=1243{,}7$~€ y el gap sube al $32{,}0\,\%$: la relajación casi no paga los costes fijos ($y_1=0{,}008$).
\end{enumerate}

\noindent\textbf{(e)}
\begin{enumerate}[label=\arabic*.]
  \item «$y_2=1\Rightarrow x_1\le 60$»: $x_1\le 60+M(1-y_2)$. Como $x_1\le 80$ siempre, la menor constante válida es $M=80-60=20$: $x_1\le 80-20\,y_2$.
  \item Solo cambian los patrones con $y_2=1$ y $x_1>60$. $\{1,2\}$: $x_1=60$, $x_2=40$: $550+720+640=1910$~€. $\{1,2,3\}$ no cambia ($x_1=55$).
  \item Nuevo óptimo: $\{2,3\}$ con $x=(0;\ 40;\ 60)$~t y $1880$~€.
\end{enumerate}
\resultado (b) $M=(80;\ 50;\ 60)$~t. (c) Hornos 1 y~2, $x=(80;\ 20;\ 0)$~t, $1830$~€. (d) $z_{\mathrm{LP}}=1723$~€ y gap del $5{,}83\,\%$ con las $M$ mínimas; $1244$~€ y $32{,}0\,\%$ con $M=10^{4}$: una $M$ grande no cambia el óptimo entero, pero destruye la cota y agranda el árbol. (e) $x_1\le 80-20y_2$; óptimo $\{2,3\}$, $x=(0;\ 40;\ 60)$~t, $1880$~€.
\comprobacion la cota es coherente: $1723{,}3\le 1830$ (proposición~2.10). En el nuevo óptimo se cumplen $L_2=20\le 40\le 50$ y $L_3=25\le 60\le 60$, y la igualdad $40+60=100$. Un PLEM con HiGHS da los mismos valores.

[Criterios: (a) 1~p por la función objetivo, 1~p por el enlace y 1~p por el lote mínimo y la igualdad. (b) 1~p por los valores y 1~p por la justificación. (c) 0{,}5~p por los patrones infactibles y 2{,}5~p por los cuatro patrones factibles y la conclusión. (d) 1~p por cada relajación con su gap. (e) 1~p por la restricción con $M=20$ y 1~p por el nuevo óptimo.]

Repasar: M24, §2.1, §2.3 y §2.4.
\end{solucion}

% ------------------------------------------------------------- B.2
\begin{solucion}{eje:bo}
\datos núcleo exponencial cuadrático con $\sigma_f=\ell=1$; $X=(0,2)$, $y=(0{,}5;\ -0{,}8)$; $\fmejor=-0{,}8$.
\planteamiento posterior del proceso gaussiano (teorema~3.5 de M25), mejora esperada en forma cerrada (teorema~4.4), probabilidad de mejora (definición~4.2) y cota inferior de confianza (definición~4.7).
\pasos

\noindent\textbf{(a)}
\begin{enumerate}[label=\arabic*.]
  \item $\rho=k(0,2)=e^{-2}=0{,}1353$; $K=\begin{psmallmatrix}1&\rho\\\rho&1\end{psmallmatrix}$, $\det K=1-\rho^{2}=0{,}9817$ y $K^{-1}=\frac{1}{0{,}9817}\begin{psmallmatrix}1&-0{,}1353\\-0{,}1353&1\end{psmallmatrix}=\begin{psmallmatrix}1{,}0187&-0{,}1379\\-0{,}1379&1{,}0187\end{psmallmatrix}$.
  \item $\alpha=K^{-1}y=\frac{1}{0{,}9817}(0{,}5+0{,}1353\cdot 0{,}8;\ -0{,}8-0{,}1353\cdot 0{,}5)=(0{,}6196;\ -0{,}8839)$.
  \item En $x=1$: $k_\ast=(e^{-1/2},e^{-1/2})=(0{,}6065;\ 0{,}6065)$. $\mu(1)=0{,}6065\,(0{,}6196-0{,}8839)=-0{,}1603$.
  \item Como $k_\ast\propto(1,1)$ y $K^{-1}(1,1)\T=(1,1)\T/(1+\rho)$: $\sigma^{2}(1)=1-2\cdot 0{,}6065^{2}/1{,}1353=1-0{,}6481=0{,}3519$; $\sigma(1)=0{,}5932$.
\end{enumerate}

\noindent\textbf{(b)}
\begin{enumerate}[label=\arabic*.]
  \item $k_\ast(3)=(e^{-9/2},e^{-1/2})=(0{,}01111;\ 0{,}6065)$.
  \item $\mu(3)=0{,}01111\cdot 0{,}6196+0{,}6065\cdot(-0{,}8839)=0{,}00688-0{,}5361=-0{,}5292$.
  \item $k_\ast\T K^{-1}k_\ast=(0{,}01111^{2}-2\cdot 0{,}1353\cdot 0{,}01111\cdot 0{,}6065+0{,}6065^{2})/0{,}9817=0{,}3661/0{,}9817=0{,}3730$; $\sigma^{2}(3)=0{,}6270$, $\sigma(3)=0{,}7918$.
\end{enumerate}

\noindent\textbf{(c)}
\begin{enumerate}[label=\arabic*.]
  \item $x=1$: $z=(-0{,}8+0{,}1603)/0{,}5932=-1{,}078$; $\Phi(z)=0{,}1404$, $\varphi(z)=0{,}2231$ (interpolando en el \cref{tab:normal}). $\PI=0{,}140$; $\EI=-0{,}6397\cdot 0{,}1404+0{,}5932\cdot 0{,}2231=-0{,}0898+0{,}1324=0{,}0425$.
  \item $x=3$: $z=(-0{,}8+0{,}5292)/0{,}7918=-0{,}3420$; $\Phi(z)=0{,}3662$, $\varphi(z)=0{,}3763$. $\PI=0{,}366$; $\EI=-0{,}2708\cdot 0{,}3662+0{,}7918\cdot 0{,}3763=-0{,}0992+0{,}2980=0{,}199$.
  \item $\LCB(1)=-0{,}1603-2\cdot 0{,}5932=-1{,}347$; $\LCB(3)=-0{,}5292-2\cdot 0{,}7918=-2{,}113$.
  \item Los tres criterios eligen $x=3$. No hace falta calcular para saberlo: $x=3$ tiene a la vez menor media y mayor desviación, y la EI decrece con $\mu$ y crece con $\sigma$ (proposición~4.5 de M25). El punto 3 combina explotación (la tendencia apunta a valores más bajos más allá del mejor ensayo) y exploración (está lejos de los datos); $x=1$, entre los dos ensayos, está bien acotado y es peor que $\fmejor$ con probabilidad alta.
\end{enumerate}

\noindent\textbf{(d)}
\begin{enumerate}[label=\arabic*.]
  \item $K_y=K+0{,}1\,I$, $\det K_y=1{,}1^{2}-0{,}1353^{2}=1{,}1917$; $k_\ast(2)=(0{,}1353;\ 1)$.
  \item $k_\ast\T K_y^{-1}k_\ast=(1{,}1\cdot 0{,}01832-2\cdot 0{,}1353\cdot 0{,}1353+1{,}1)/1{,}1917=1{,}0835/1{,}1917=0{,}9092$; $\sigma^{2}_{f}(2)=0{,}0908$.
  \item Varianza de una nueva observación: $0{,}0908+0{,}1=0{,}191$.
  \item En la EI se usa la latente, $0{,}0908$: la mejora se refiere al valor verdadero $f(x)$, no a una medida ruidosa. Con la predictiva, la incertidumbre de un punto ensayado nunca baja de $\sigma_\varepsilon$ y el algoritmo repite ensayos (observación~4.8 de M25).
\end{enumerate}

\noindent\textbf{(e)}
\begin{enumerate}[label=\arabic*.]
  \item Con $\ell=0{,}3$: $k(0,1)=k(2,1)=e^{-0{,}5/0{,}09}=e^{-5{,}556}=0{,}00387$ y $k(0,2)=e^{-22{,}2}\approx 0$, de modo que $K\approx I$ y $\alpha\approx y$.
  \item $\mu(1)\approx 0{,}00387\,(0{,}5-0{,}8)=-0{,}0012\approx 0$; $\sigma^{2}(1)\approx 1-2\cdot 0{,}00387^{2}\approx 1$.
  \item $z\approx-0{,}8$: $\EI(1)\approx-0{,}8\cdot 0{,}2119+0{,}2897=0{,}120$. En $x=3$ sale prácticamente lo mismo ($0{,}121$).
  \item Con una longitud de escala corta los datos no informan entre ensayos: toda la zona no ensayada vuelve a la previa y la EI es casi uniforme lejos de los datos, de modo que el algoritmo explora casi al azar. Por eso $\ell$ se estima por verosimilitud marginal (M25, §3.4) y se vigila que no colapse.
\end{enumerate}
\resultado (a) $\alpha=(0{,}620;\ -0{,}884)$, $\mu(1)=-0{,}160$, $\sigma^{2}(1)=0{,}352$. (b) $\mu(3)=-0{,}529$, $\sigma^{2}(3)=0{,}627$. (c) $\PI=(0{,}140;\ 0{,}366)$, $\EI=(0{,}0425;\ 0{,}199)$, $\LCB=(-1{,}35;\ -2{,}11)$; los tres criterios eligen $x=3$. (d) $0{,}0908$ (latente) y $0{,}191$ (observación); se usa la latente. (e) $\mu(1)\approx 0$, $\sigma^{2}(1)\approx 1$, $\EI(1)\approx 0{,}120$.
\comprobacion la varianza nunca supera la previa ($0{,}352<1$, $0{,}627<1$). Ambos candidatos distan una unidad de su ensayo más próximo, pero $x=1$ tiene \emph{dos} ensayos a esa distancia y $x=3$ solo uno, de ahí que su varianza sea menor. Una búsqueda en rejilla de la EI con $\ell=1$ da el máximo en $x\approx 2{,}95$ con $\EI=0{,}199$.

[Criterios: (a) 1~p por $K^{-1}$ y $\alpha$, 1~p por $\mu(1)$ y 1~p por $\sigma^{2}(1)$. (b) 1~p por cada cifra. (c) 1~p por las EI, 1~p por PI y LCB y 1~p por la explicación. (d) 1~p por los cálculos y 1~p por la elección razonada. (e) 1~p por las cifras y 1~p por la interpretación.]

Repasar: M25, §3.3, §3.4, §4.2, §4.3, §4.4 y §4.5.
\end{solucion}

% ------------------------------------------------------------- B.3
\begin{solucion}{eje:mpc}
\datos $a=0{,}9$, $b=0{,}5$; $Q=1$, $R=0{,}5$, $P_f=1$, $N=2$, $x_0=2$; en~(c), $q=2$, $r=1$, $\gamma=0{,}95$.
\planteamiento (a) recursión de Riccati del teorema~3.2 de M26; (b) forma condensada y QP de la proposición~3.6, resuelto con las condiciones de Karush--Kuhn--Tucker como en el ejemplo~3.7; (c) ecuación de Bellman (M26, §3.1 y §4.1) con una función de valor afín ensayada.
\pasos

\noindent\textbf{(a)}
\begin{enumerate}[label=\arabic*.]
  \item $P_2=1$. $K_1=bP_2a/(R+b^{2}P_2)=0{,}45/0{,}75=0{,}6$; $P_1=Q+a^{2}P_2-aP_2bK_1=1+0{,}81-0{,}27=1{,}54$.
  \item $K_0=bP_1a/(R+b^{2}P_1)=0{,}693/0{,}885=0{,}7831$; $P_0=1+0{,}81\cdot 1{,}54-0{,}9\cdot 1{,}54\cdot 0{,}5\cdot 0{,}7831=1+1{,}2474-0{,}5426=1{,}7047$.
  \item $u_0=-0{,}7831\cdot 2=-1{,}566$; $x_1=1{,}8-0{,}7831=1{,}017$; $u_1=-0{,}6\cdot 1{,}017=-0{,}6102$; $x_2=0{,}9153-0{,}3051=0{,}6102$.
  \item Coste: $4+0{,}5\cdot 2{,}453+1{,}034+0{,}5\cdot 0{,}3723+0{,}3723=6{,}819=P_0x_0^{2}=1{,}7047\cdot 4$.
\end{enumerate}

\noindent\textbf{(b)}
\begin{enumerate}[label=\arabic*.]
  \item $\Phi=(a,a^{2})\T=(0{,}9;\ 0{,}81)\T$, $\Gamma=\begin{psmallmatrix}b&0\\ab&b\end{psmallmatrix}=\begin{psmallmatrix}0{,}5&0\\0{,}45&0{,}5\end{psmallmatrix}$, $\bar Q=\operatorname{diag}(1,1)$, $\bar R=\operatorname{diag}(0{,}5;\ 0{,}5)$.
  \item $H=\Gamma\T\bar Q\Gamma+\bar R=\begin{psmallmatrix}0{,}9525&0{,}225\\0{,}225&0{,}75\end{psmallmatrix}$; $F=\Gamma\T\bar Q\Phi=(0{,}8145;\ 0{,}405)$; constante $x_0^{2}(1+0{,}81+0{,}6561)=9{,}864$.
  \item El óptimo libre tiene $u_0=-1{,}566<-1$: viola la cota. Se ensaya $u_0=-1$ activo y se minimiza en $u_1$: $\partial J/\partial u_1=2(0{,}225u_0+0{,}75u_1)+2\cdot 2\cdot 0{,}405=0$ da $u_1=(0{,}225-0{,}81)/0{,}75=-0{,}78$, dentro de la cota.
  \item $J=7{,}103$ (frente a $6{,}819$ sin restricción). En ese punto $\partial J/\partial u_0=2(0{,}9525\cdot(-1)+0{,}225\cdot(-0{,}78))+2\cdot 2\cdot 0{,}8145=1{,}002>0$: subir $u_0$ empeora y bajarlo no está permitido; el multiplicador de $u_0\ge-1$ es $1{,}002$.
  \item Se aplica solo $u_0=-1$~kW; en el instante siguiente se mide $x_1$ y se vuelve a resolver (horizonte deslizante, M26, §3.2).
\end{enumerate}

\noindent\textbf{(c)}
\begin{enumerate}[label=\arabic*.]
  \item Bellman: $V(z)=\min_u\bigl\{qz+ru^{2}+\gamma V(az+bu)\bigr\}$.
  \item Con $V(z)=pz+s$: $pz+s=qz+\gamma paz+\gamma s+\min_u\{ru^{2}+\gamma pbu\}$.
  \item El mínimo en $u$ no depende de $z$: $u^{\ast}=-\gamma pb/(2r)$, con valor $-(\gamma pb)^{2}/(4r)$.
  \item Identificando coeficientes: $p=q+\gamma ap$, luego $p=q/(1-\gamma a)=2/0{,}145=13{,}79$; y $s(1-\gamma)=-(\gamma pb)^{2}/(4r)$.
  \item $\gamma pb=0{,}95\cdot 13{,}79\cdot 0{,}5=6{,}552$; $u^{\ast}=-3{,}276$~kW; $s=-6{,}552^{2}/(4\cdot 0{,}05)=-214{,}6$.
\end{enumerate}

\noindent\textbf{(d)}
\begin{enumerate}[label=\arabic*.]
  \item Si $V(z)=pz+s$, el término en $z$ de la ecuación de Bellman no depende de $u$, de modo que la minimización en $u\in[-2,2]$ sigue sin depender de $z$ y la identificación $p=q+\gamma ap$ no cambia: la forma afín se conserva con la misma $p$.
  \item La función $ru^{2}+\gamma pbu$ es convexa con mínimo libre en $-3{,}276$, fuera del intervalo: $u^{\ast}=-2$~kW (saturada).
  \item $s(1-\gamma)=r\cdot 4+6{,}552\cdot(-2)=-9{,}103$, $s=-182{,}1$.
\end{enumerate}
\resultado (a) $K_1=0{,}600$, $P_1=1{,}54$, $K_0=0{,}783$, $P_0=1{,}70$; $u=(-1{,}57;\ -0{,}610)$~kW y $J=6{,}82=P_0x_0^{2}$. (b) $u=(-1;\ -0{,}780)$~kW, $J=7{,}10$, multiplicador $1{,}00$; se aplica $u_0=-1$~kW. (c) $V(z)=13{,}8\,z-215$ y $u^{\ast}=-3{,}28$~kW, constante. (d) $u^{\ast}=-2$~kW y $V(z)=13{,}8\,z-182$.
\comprobacion en~(c), simulando la política constante desde $z=1$ durante 3000 pasos se obtiene $J=-200{,}83=13{,}79-214{,}63$, y perturbar $u$ en $\pm0{,}1$ empeora el coste en $0{,}2$ en ambos sentidos. La condición $\gamma a=0{,}855<1$ garantiza que la serie converge. La acción no depende del estado porque el coste marginal de tener una unidad más de $x$ es el mismo en cualquier estado: es lo que dice la pendiente constante $p$.

[Criterios: (a) 1~p por $K_1,P_1$, 1~p por $K_0,P_0$ y 1~p por la secuencia y la comprobación. (b) 1~p por $H$ y $F$, 1~p por la solución con la cota activa y 1~p por el multiplicador y el principio de horizonte deslizante. (c) 1~p por la ecuación de Bellman, 1~p por $u^{\ast}$, 1~p por $p$ y 1~p por $s$. (d) 1~p por el argumento y 1~p por $u^{\ast}$ y $s$.]

Repasar: M26, §3.1, §3.2, §3.3 y §4.1.
\end{solucion}

% ------------------------------------------------------------- B.4
\begin{solucion}{eje:simp}
\datos $k_0=1000$~N/mm, $p=3$, $x=(1;\ 0{,}8;\ 0{,}5)$, $F=(0,100,100)$~N.
\planteamiento la barra en serie es isostática; el gradiente de la flexibilidad se obtiene con el corolario~2.5 de M27 y el de $J=e_2\T U$ con el teorema~2.4 (método adjunto); la actualización es la de la definición~2.7, y con $p=1$ el problema es convexo (teorema~2.3), de modo que las condiciones de Karush--Kuhn--Tucker son suficientes (M24, §4.2).
\pasos

\noindent\textbf{(a)}
\begin{enumerate}[label=\arabic*.]
  \item $k=(1000;\ 512;\ 125)$~N/mm. $K=\begin{psmallmatrix}k_1+k_2&-k_2&0\\-k_2&k_2+k_3&-k_3\\0&-k_3&k_3\end{psmallmatrix}=\begin{psmallmatrix}1512&-512&0\\-512&637&-125\\0&-125&125\end{psmallmatrix}$.
  \item Esfuerzos axiles (equilibrio de cada tramo): $N=(200;\ 200;\ 100)$~N. Alargamientos $\delta_e=N_e/k_e=(0{,}2;\ 0{,}3906;\ 0{,}8)$~mm.
  \item $U=(0{,}2;\ 0{,}5906;\ 1{,}3906)$~mm; $c=100\cdot 0{,}5906+100\cdot 1{,}3906=198{,}1$~N·mm.
\end{enumerate}

\noindent\textbf{(b)}
\begin{enumerate}[label=\arabic*.]
  \item $g=F\T U$ no depende de $x$ y $(\partial g/\partial U)\T=F$; el adjunto resuelve $K\lambda=F$, luego $\lambda=U$.
  \item $\partial c/\partial x_e=-U\T(\partial K/\partial x_e)U=-p\,x_e^{p-1}k_0\,\delta_e^{2}$: $-3\cdot 1\cdot 1000\cdot 0{,}04=-120$; $-3\cdot 0{,}64\cdot 1000\cdot 0{,}1526=-293{,}0$; $-3\cdot 0{,}25\cdot 1000\cdot 0{,}64=-480$~N·mm.
  \item Comprobación directa: $c=\sum_eN_e^{2}/(k_0x_e^{3})$, con $N$ fijo; $\partial c/\partial x_3=-3N_3^{2}/(k_0x_3^{4})=-3\cdot 10^{4}/(1000\cdot 0{,}0625)=-480$.
\end{enumerate}

\noindent\textbf{(c)}
\begin{enumerate}[label=\arabic*.]
  \item Adjunto: $K\lambda=e_2$, es decir, el desplazamiento producido por una carga unidad en el nodo~2: $\lambda=(1/k_1;\ 1/k_1+1/k_2;\ 1/k_1+1/k_2)=(0{,}001;\ 0{,}002953;\ 0{,}002953)$~mm/N, con alargamientos $\delta^{\lambda}=(0{,}001;\ 0{,}001953;\ 0)$.
  \item $\partial J/\partial x_e=-\lambda\T(\partial K/\partial x_e)U=-p\,x_e^{p-1}k_0\,\delta^{\lambda}_e\,\delta_e$: $-3\cdot 1000\cdot 0{,}001\cdot 0{,}2=-0{,}6$; $-3\cdot 0{,}64\cdot 1000\cdot 0{,}001953\cdot 0{,}3906=-1{,}465$; $0$~mm.
  \item El elemento~3 está aguas abajo del nodo~2: su rigidez no cambia $u_2$, porque la barra es isostática. Comprobación: $u_2=200/k_1+200/k_2$ no depende de $x_3$ y $\partial u_2/\partial x_1=-200\cdot 3000/10^{6}=-0{,}6$.
\end{enumerate}

\noindent\textbf{(d)}
\begin{enumerate}[label=\arabic*.]
  \item $B_e=(-\partial c/\partial x_e)/\Lambda$ con $\partial V/\partial x_e=1$. Con $\Lambda=300$: $B=(0{,}4;\ 0{,}9766;\ 1{,}6)$, $\sqrt B=(0{,}6325;\ 0{,}9882;\ 1{,}2649)$.
  \item $x\sqrt B=(0{,}6325;\ 0{,}7906;\ 0{,}6325)$. Los límites de movimiento son $[0{,}8;\,1]$, $[0{,}6;\,1]$ y $[0{,}3;\,0{,}7]$, de modo que $x^{\text{nuevo}}=(0{,}8;\ 0{,}7906;\ 0{,}6325)$, con suma $2{,}223<2{,}3$.
  \item Falta material: $\Lambda$ (el precio del material) es demasiado alto y debe bajar (proposición~2.8: el volumen decrece con $\Lambda$).
  \item Con $\Lambda=270$: $\sqrt{120/270}=0{,}6667$ (se recorta a $0{,}8$), $\sqrt{292{,}97/270}=1{,}0417$ y $\sqrt{480/270}=1{,}3333$; $x^{\text{nuevo}}=(0{,}8;\ 0{,}8333;\ 0{,}6667)$, suma $2{,}3$.
\end{enumerate}

\noindent\textbf{(e)}
\begin{enumerate}[label=\arabic*.]
  \item Con $p=1$: $c(x)=\sum_eN_e^{2}/(k_0x_e)$, convexa en $x>0$. Lagrangiana $c+\Lambda(\sum_ex_e-2{,}3)$; estacionariedad: $N_e^{2}/(k_0x_e^{2})=\Lambda$, luego $x_e\propto N_e$.
  \item $x_e=2{,}3\,N_e/\sum N=2{,}3\cdot(200,200,100)/500=(0{,}92;\ 0{,}92;\ 0{,}46)$, todas $\le 1$; $\Lambda=(500/2{,}3)^{2}/1000=47{,}26$~N·mm por unidad de volumen.
  \item $c^{\ast}=\frac{1}{1000}\bigl(\frac{4\cdot 10^{4}}{0{,}92}\cdot 2+\frac{10^{4}}{0{,}46}\bigr)=108{,}7$~N·mm, frente a $c(x_0)=\frac{1}{1000}\bigl(\frac{4\cdot 10^{4}}{1}+\frac{4\cdot 10^{4}}{0{,}8}+\frac{10^{4}}{0{,}5}\bigr)=110{,}0$~N·mm.
\end{enumerate}
\resultado (a) Desplazamientos $U=(0{,}200;\ 0{,}591;\ 1{,}39)$~mm y flexibilidad $c=198$~N·mm. (b) Gradiente $\nabla c=(-120;\ -293;\ -480)$~N·mm. (c) Gradiente $\nabla J=(-0{,}600;\ -1{,}46;\ 0)$~mm. (d) Con $\Lambda=300$ resulta $x=(0{,}800;\ 0{,}791;\ 0{,}633)$ y un volumen de $2{,}22$, de modo que hay que bajar $\Lambda$; con $\Lambda=270$ resulta $x=(0{,}800;\ 0{,}833;\ 0{,}667)$ y un volumen de $2{,}30$. (e) $x^{\ast}=(0{,}920;\ 0{,}920;\ 0{,}460)$ y $c^{\ast}=109$~N·mm, un $1{,}2\,\%$ menos que el diseño inicial con $p=1$.
\comprobacion las diferencias finitas con paso $10^{-6}$ reproducen los tres gradientes de~(b) y de~(c) con cinco cifras. En~(e), en el óptimo $N_e^{2}/(k_0x_e^{2})$ vale $47{,}26$ en los tres elementos, como exige la estacionariedad, y el signo negativo de todas las sensibilidades de~(b) confirma que añadir material nunca aumenta la flexibilidad.

[Criterios: (a) 1~p por $K$ y 1~p por $U$ y $c$. (b) 1~p por el argumento autoadjunto, 1~p por las tres derivadas y 1~p por la comprobación. (c) 1~p por el problema adjunto, 1~p por las derivadas y 1~p por la interpretación. (d) 1~p por $x^{\text{nuevo}}$ y el sentido de $\Lambda$ y 1~p por la verificación con $\Lambda=270$. (e) 1~p por las condiciones KKT y 1~p por el óptimo y la comparación.]

Repasar: M27, §2.1, §2.2, §2.3 y §2.4; M24, §4.2.
\end{solucion}

% ------------------------------------------------------------- B.5
\begin{solucion}{eje:valor}
\datos $n=16$, $\bar d=1200$~€/semana, $s_d=1600$~€/semana; $\sigma_A=\sigma_B=3000$~€, $r=0{,}86$; percentiles \emph{bootstrap} 560 y 1890~€/semana; $I=60\,000$~€, $O=12\,000$~€/año, $N=5$, $\rho=0{,}08$.
\planteamiento intervalo $t$ pareado y ventaja del emparejamiento (proposición~5.3 de M28), intervalos \emph{bootstrap} (definición~5.4), VAN con anualidad (proposición~5.2) y tamaño de un piloto (proposición~5.5).
\pasos

\noindent\textbf{(a)}
\begin{enumerate}[label=\arabic*.]
  \item Error típico $s_d/\sqrt n=1600/4=400$; semiamplitud $2{,}131\cdot 400=852{,}4$.
  \item Intervalo: $[347{,}6;\ 2052{,}4]$~€/semana, es decir, $[18\,075;\ 106\,725]$~€/año.
  \item Hipótesis: diferencias independientes entre semanas (o autocorrelación despreciable, que debe comprobarse) y aproximadamente normales, o $n$ suficiente; el gemelo debe estar validado, porque el intervalo mide la variabilidad entre semanas y no el error del modelo.
\end{enumerate}

\noindent\textbf{(b)}
\begin{enumerate}[label=\arabic*.]
  \item $\sigma_D=\sqrt{3000^{2}+3000^{2}-2\cdot 0{,}86\cdot 3000^{2}}=\sqrt{2{,}52\cdot 10^{6}}=1587$~€, coherente con $s_d$.
  \item Sin emparejar: $\sqrt{2\cdot 3000^{2}}=4243$~€.
  \item Cociente de varianzas: $1/(1-r)=7{,}14$: harían falta unas siete veces más semanas.
\end{enumerate}

\noindent\textbf{(c)}
\begin{enumerate}[label=\arabic*.]
  \item Percentil: $[560;\ 1890]$~€/semana.
  \item Básico: $[2\bar d-1890;\ 2\bar d-560]=[510;\ 1840]$~€/semana.
  \item La distribución de las diferencias es asimétrica: la mayoría de semanas el ahorro es moderado y, de vez en cuando, el optimizador evita una avería cara. El \emph{bootstrap} recoge esa asimetría sin suponer normalidad; ambos métodos suponen semanas independientes.
\end{enumerate}

\noindent\textbf{(d)}
\begin{enumerate}[label=\arabic*.]
  \item $a_{5;\,0{,}08}=(1-1{,}08^{-5})/0{,}08=3{,}9927$.
  \item Con la media: $b=52\cdot 1200=62\,400$~€/año; $\VAN=(62\,400-12\,000)\cdot 3{,}9927-60\,000=141\,233$~€.
  \item Con el extremo inferior: $b=18\,075$; $\VAN=(18\,075-12\,000)\cdot 3{,}9927-60\,000=-35\,744$~€.
  \item Umbral: $\VAN=0\iff b=O+I/a=12\,000+15\,027=27\,027$~€/año, es decir, $519{,}8$~€/semana.
  \item Como $347{,}6<519{,}8$, el intervalo del 95\,\% incluye mejoras que no recuperan la inversión: no puede afirmarse con esa confianza que el proyecto sea rentable, aunque el VAN esperado sea claramente positivo. (El contraste unilateral de $H_0\colon\delta\le 519{,}8$ da $t=(1200-519{,}8)/400=1{,}70<1{,}753$.)
\end{enumerate}

\noindent\textbf{(e)}
\begin{enumerate}[label=\arabic*.]
  \item $z_{0{,}975}+z_{0{,}80}=1{,}960+0{,}8416=2{,}802$.
  \item $n\approx(2{,}802\cdot 1600/519{,}8)^{2}=(8{,}624)^{2}=74{,}4$: unas 75 semanas.
\end{enumerate}
\resultado (a) $[348;\ 2050]$~€/semana, es decir, $[18\,100;\ 107\,000]$~€/año. (b) $\sigma_D=1590$~€ frente a $4240$~€; unas $7{,}1$ veces más semanas. (c) Percentil $[560;\ 1890]$ y básico $[510;\ 1840]$~€/semana. (d) $\VAN=141\,000$~€ con la media y $-35\,700$~€ con el extremo inferior; la rentabilidad nula exige $520$~€/semana; no puede afirmarse con un 95\,\% de confianza. (e) Unas 75 semanas.
\comprobacion la recuperación con la mejora media es $60\,000/50\,400=1{,}19$ años y la tasa interna de rentabilidad, del $79\,\%$; con $\delta=1200$ bastarían $n\approx 14$ semanas, lo que explica que 16 semanas detecten la mejora pero no basten para demostrar la rentabilidad. Las conclusiones serían las de M28, §5.3: completar con más líneas, un periodo en sombra o indicadores con menos ruido.

[Criterios: (a) 1~p por la semiamplitud, 1~p por el intervalo en las dos escalas y 1~p por las hipótesis. (b) 1~p por las dos desviaciones y 1~p por el factor. (c) 1~p por los dos intervalos y 1~p por la explicación. (d) 1~p por los dos VAN, 1~p por el umbral y 1~p por la conclusión. (e) 2~p, o 1~p si se usa $\delta=1200$ sin justificarlo.]

Repasar: M28, §5.1, §5.2 y §5.3.
\end{solucion}

% =====================================================================
\chapter{Orientaciones de corrección del proyecto}
\label{ap:correccion}
% =====================================================================

Este apéndice está pensado para el profesor y para el propio estudiante, que debería aplicarlo a su trabajo antes de entregarlo. La corrección no busca una cifra de mejora concreta, sino comprobar que la cifra que se presenta es cierta, que se ha medido contra referencias honradas y que el sistema no viola nada que no pueda violar.

\section{Orden de la corrección}

\begin{enumerate}[label=\arabic*.]
  \item \emph{Reproducibilidad}. Se clona el repositorio en una máquina limpia y se ejecuta la orden única en el modo reducido documentado. Se comprueba que la huella del conjunto de prueba coincide con la registrada en la semana~1 y que las cifras principales del informe se reproducen con una diferencia relativa menor del 0{,}5\,\%. Si no se reproducen, el criterio «Código y demostración» se califica como insuficiente y las cifras no reproducidas no se tienen en cuenta en el resto de la rúbrica.
  \item \emph{Criterios obligatorios}. Se recorren los criterios A, B o C de la opción, con la evidencia que señala el informe (tablas, registros de MLflow, salidas de las pruebas automáticas).
  \item \emph{Rúbrica}. Se puntúa cada criterio con la tabla de la opción. En caso de duda entre dos niveles, decide la calidad de la justificación escrita, no el resultado numérico.
  \item \emph{Defensa}. Las preguntas de la demostración sirven para confirmar la autoría y detectar deuda técnica oculta~\cite{sculley}; también pueden mover la nota de cada criterio un nivel en cualquier sentido.
\end{enumerate}

\section{Comprobación de la evaluación honesta}

\begin{table}[H]
\centering\small
\caption{Cómo se comprueban las reglas R1 a R7 de la \cref{sec:proy-comun}.}\label{tab:honesta}
\begin{tabularx}{\textwidth}{@{}c>{\raggedright\arraybackslash}X>{\raggedright\arraybackslash}X@{}}
\toprule
Regla & Qué se mira & Señal de alarma \\
\midrule
R1 & El generador acepta la semilla y reproduce el escenario nominal; los parámetros coinciden con los del enunciado & Parámetros cambiados para facilitar la mejora sin justificarlo \\
R2 & Fecha y huella del conjunto de prueba; registros de experimentos & Ejecuciones sobre la prueba antes de la semana~7; ajuste de hiperparámetros con las métricas de prueba \\
R3 & Tabla con línea base, alternativa sencilla, optimizador y techo & Solo se compara con la práctica actual; techo ausente \\
R4 & Mismas semanas, trayectorias o días y mismas semillas en todas las políticas & Políticas evaluadas en periodos o con ruidos distintos \\
R5 & Intervalo $t$ pareado y \emph{bootstrap}; comprobación de autocorrelación & Desviación típica de cada política en lugar de la de la diferencia \\
R6 & Prueba automática que viola a propósito una restricción y comprueba que la salvaguarda actúa & Restricciones duras solo dentro del modelo de optimización \\
R7 & Declaración del uso de asistentes y de cómo se verificó lo generado & Código que el estudiante no sabe explicar en la defensa \\
\bottomrule
\end{tabularx}
\end{table}

\section{Errores frecuentes y su efecto en la nota}

\begin{table}[H]
\centering\small
\caption{Errores frecuentes en los proyectos y criterio afectado.}\label{tab:errores-proyecto}
\begin{tabularx}{\textwidth}{@{}>{\raggedright\arraybackslash}X>{\raggedright\arraybackslash}p{3.8cm}>{\raggedright\arraybackslash}p{3.0cm}@{}}
\toprule
Error & Remisión & Criterio afectado \\
\midrule
Big-M de $10^{6}$ «por si acaso»; \emph{gap} sin informar & M24, §2.3 y §2.4 & formulación \\
Previsión evaluada solo por su error, no por la decisión que induce & M24, §6.1 & componente de IA \\
Varianza predictiva en lugar de latente en la mejora esperada & M25, §4.5 & componente de IA \\
Episodios truncados tratados como terminales en el RL & M26, §4.4 & control \\
Comparar un MPC con previsiones con un PID sin ellas como si fuera una comparación de métodos & M26, §5.1 & evaluación \\
Restricción de rendimiento sobre la media del tiempo de fotograma & M27, §5.1 & selección con presupuesto \\
LOD sin validar: caras degeneradas, orientación incoherente & M27, §6.2 & cadena de LOD \\
Censura ignorada: vidas interrumpidas tratadas como fallos o descartadas & M28, §4.5 & componente de IA \\
Sesgo del ganador: informar la puntuación de validación del mejor candidato & M28, §2.2 & evaluación \\
VAN calculado frente a la práctica actual y atribuido entero a la IA & M28, §5.1 & valor \\
\bottomrule
\end{tabularx}
\end{table}

\section{Cifras de control}

Las siguientes magnitudes se deducen del escenario nominal con los resultados de los apuntes y permiten detectar generadores o gemelos mal construidos. No son resultados esperados del optimizador.
\begin{itemize}[label=\textendash]
  \item \emph{Opción (a)}. En régimen permanente, mantener el horno a 530~°C con 20~°C de ambiente exige $(530-20)/R=510$~kW, y durante la primera hora de una carga fría unos 150~kW más; con los dos centros de mecanizado en marcha (300~kW), un calentamiento a plena potencia (800~kW) supera la potencia contratada, de modo que esa restricción debe estar activa en algunas horas del plan. La constante de tiempo del horno es $RC=3$~h.
  \item \emph{Opción (b)}. Todas las texturas en BC7 y BC5 a resolución completa ocupan $40\cdot 3\cdot 4096^{2}\cdot\tfrac43$~bytes $\approx 2{,}68$~GB, por encima del presupuesto de 1{,}5~GB. Con $H=1080$ y $60^{\circ}$ de campo de visión, un error de 1~cm alcanza un píxel a $9{,}35$~m (M27, proposición~5.5). A 60~Hz el presupuesto es de $16{,}7$~ms (M27, definición~5.1).
  \item \emph{Opción (c)}. Con la demanda media, la capacidad exige unas $650\cdot 2{,}5/70=23{,}2$ furgonetas, y los lunes, $27{,}9$. Sin covariables, la política por edad de la proposición~4.13 de M28 con $\beta=2{,}5$, $\theta=15\,000$~km, $c_p=350$~€ y $c_f=1800$~€ tiene su óptimo en $T^{\ast}\approx 7300$~km, con un coste de unos 81{,}9~€ por cada 1000~km, frente a 87{,}3~€ con la revisión cada 10\,000~km y 135~€ sin mantenimiento preventivo. La probabilidad de fallar antes de 10\,000~km es $1-e^{-(2/3)^{2{,}5}}=0{,}30$.
\end{itemize}

\section{Preguntas de la defensa}

\begin{enumerate}[label=\arabic*.]
  \item \emph{Comunes}. ¿Qué parte de la mejora frente a la práctica actual se debe a la IA y qué parte a la alternativa sencilla? ¿Qué fracción del margen hasta el techo se captura y qué haría falta para capturar más? ¿Qué semana del conjunto de prueba es la peor y por qué? ¿Qué ocurre si se duplica el ruido del generador?
  \item \emph{Opción (a)}. ¿Cuál es el precio sombra de la potencia contratada y qué decisión de compra de potencia justificaría? ¿Por qué no basta planificar con el precio medio previsto? Muestre el filtro de seguridad actuando con un sensor congelado.
  \item \emph{Opción (b)}. ¿Qué cota da la relajación lagrangiana en un fotograma y qué hueco de dualidad tiene? ¿Cómo se ajustaron las ganancias del controlador de $\lambda$ y qué pasa si la GPU es un 30\,\% más lenta? ¿Qué comprobación geométrica falló más a menudo durante el desarrollo?
  \item \emph{Opción (c)}. ¿Cómo trata el modelo las vidas censuradas? ¿Cuánto cambia $T^{\ast}$ si $c_f$ se duplica? ¿Coincide la estimación por ponderación inversa de la propensión con el gemelo, y por qué podría no hacerlo?
\end{enumerate}

% =====================================================================
\chapter{Plantilla del documento de diseño}
\label{ap:plantilla}
% =====================================================================

La plantilla desarrolla la de M28, §6.2, para el proyecto final, siguiendo la revisión preliminar de diseño del curso 16.842 del MIT~\cite{mit16842} y el planteamiento de sistemas con aprendizaje automático de CS329S~\cite{cs329s,huyen}. Se rellena antes de escribir código y se actualiza en las semanas~2, 5 y~8 (entregable~E1). Cada sección indica qué debe responder, qué evidencia numérica se exige y en qué versión debe estar completa. Las secciones vacías se marcan como «pendiente», no se eliminan.

\section{Ficha resumen}

La primera página del documento es esta ficha, con una o dos líneas por campo.

\begin{center}
\small
\begin{tabularx}{\textwidth}{@{}>{\raggedright\arraybackslash}p{4.6cm}X@{}}
\toprule
Campo & Contenido \\
\midrule
Opción y título & \\
Autor o autores y reparto & \\
Decisión, frecuencia y horizonte & \\
Objetivo (escalar, con unidades) & \\
Restricciones duras & \\
Indicadores que se vigilan & \\
Línea base / alternativa sencilla / techo & tres cifras con unidades y su origen \\
Formulación y resolvedor & \\
Componente aprendido & qué predice, con qué métrica de decisión \\
Salvaguardas y respaldo & \\
Mejora en prueba (IC del 95\,\%) & frente a la línea base y frente a la alternativa sencilla \\
VAN y $\Prob(\VAN<0)$ & \\
Tres riesgos principales (AMFE) & \\
Orden de reproducción y tiempo & \\
\bottomrule
\end{tabularx}
\end{center}

\section{Secciones del documento}

\begin{table}[H]
\centering\footnotesize
\caption{Secciones del documento de diseño, evidencia exigida y versión en que deben estar completas.}\label{tab:plantilla-secciones}
\begin{tabularx}{\textwidth}{@{}>{\raggedright\arraybackslash}p{2.6cm}>{\raggedright\arraybackslash}X>{\raggedright\arraybackslash}p{3.6cm}c@{}}
\toprule
Sección & Preguntas que debe responder & Evidencia numérica & Versión \\
\midrule
1. Contexto y decisión & ¿Qué se decide, cada cuánto, con qué horizonte y con qué información disponible en el instante de decidir? ¿Qué pasa hoy? & frecuencias, horizontes, volumen de decisiones & v0 \\
2. Partes interesadas & ¿Quién decide hoy, quién ejecuta, quién responde de los errores y quién puede parar el sistema? & --- & v0 \\
3. Objetivo y requisitos & Objetivo escalar; indicadores; restricciones duras y blandas; requisitos de latencia, disponibilidad y calidad con márgenes (\cref{tab:plantilla-requisitos}) & tabla de requisitos con umbrales y márgenes & v0 \\
4. Línea base y techo & Coste de la práctica actual, de la mejor alternativa sencilla y con información perfecta; margen máximo alcanzable & tres cifras con su procedimiento & v0 \\
5. Datos y gemelo & Generador (ecuaciones, parámetros, semillas); validación del gemelo; partición y huella de la prueba & pruebas de validación del gemelo & v0 \\
6. Formulación y método & Modelo matemático completo; familia de métodos elegida y alternativas descartadas con su motivo (mapa de M24, §1.3) & tamaño del modelo, constantes big-M justificadas & v0 \\
7. Modelos y validación & Qué se predice; validación temporal; métricas de predicción y de decisión; selección sin sesgo del ganador & métricas en validación & v1 \\
8. Salvaguardas y respaldo & Comprobaciones independientes; política de respaldo; tiempo máximo de cálculo; modos de operación & prueba automática de cada salvaguarda & v1 \\
9. Evaluación fuera de línea & Protocolo con CRN; resultados frente a las tres referencias con intervalos; ablaciones; pruebas de estrés (\cref{tab:plantilla-resultados}) & tabla de resultados & final \\
10. Piloto & Diseño de un piloto real: modo sombra, tamaño (M28, proposición~5.5), duración, criterio de aceptación fijado de antemano & tamaño del piloto & final \\
11. Despliegue y supervisión & API o herramienta, interfaz, registro de decisiones, versiones y vuelta atrás & latencias medidas & final \\
12. Monitorización & Señales, umbrales, acciones y periodo de reentrenamiento (M28, proposición~4.8) & umbrales calibrados & final \\
13. Valor y riesgo & VAN con incertidumbre, CVaR, AMFE (\cref{tab:plantilla-amfe}), encaje normativo (M28, §5.6) & VAN, $\Prob(\VAN<0)$, NPR & final \\
14. Ciclo de vida & Responsable del mantenimiento, revisiones periódicas, criterios de retirada & --- & final \\
\bottomrule
\end{tabularx}
\end{table}

\section{Tablas que se rellenan}

\begin{table}[H]
\centering\footnotesize
\caption{Plantilla de requisitos (sección~3). Se muestra una fila de ejemplo de la opción~(a).}\label{tab:plantilla-requisitos}
\begin{tabularx}{\textwidth}{@{}c>{\raggedright\arraybackslash}X>{\raggedright\arraybackslash}p{2.3cm}>{\raggedright\arraybackslash}p{1.6cm}>{\raggedright\arraybackslash}p{1.6cm}>{\raggedright\arraybackslash}p{2.6cm}@{}}
\toprule
Id. & Requisito & Métrica & Umbral & Margen de diseño & Verificación \\
\midrule
RQ1 & No superar la potencia contratada & potencia horaria máxima & 1000~kW & 950~kW & prueba automática y gemelo \\
\rule{0pt}{14pt}RQ2 & & & & & \\
\rule{0pt}{14pt}RQ3 & & & & & \\
\bottomrule
\end{tabularx}
\end{table}

\begin{table}[H]
\centering\footnotesize
\caption{Plantilla de resultados (sección~9): una fila por política, en el conjunto de prueba, con CRN y al menos 5 semillas.}\label{tab:plantilla-resultados}
\begin{tabularx}{\textwidth}{@{}>{\raggedright\arraybackslash}X cccccc@{}}
\toprule
Política & \makecell{Coste\\medio} & \makecell{Mejora frente\\a línea base} & \makecell{IC 95\,\%\\$t$ pareado} & \makecell{IC 95\,\%\\\emph{bootstrap}} & \makecell{Servicio} & \makecell{Violaciones\\duras} \\
\midrule
\rule{0pt}{13pt}Línea base & & --- & --- & --- & & \\
\rule{0pt}{13pt}Alternativa sencilla & & & & & & \\
\rule{0pt}{13pt}Optimizador con IA & & & & & & \\
\rule{0pt}{13pt}Optimizador sin componente aprendido & & & & & & \\
\rule{0pt}{13pt}Techo (información perfecta) & & & & & & \\
\bottomrule
\end{tabularx}
\end{table}

\begin{table}[H]
\centering\footnotesize
\caption{Plantilla de AMFE (sección~13), según la definición~5.8 de M28.}\label{tab:plantilla-amfe}
\begin{tabularx}{\textwidth}{@{}>{\raggedright\arraybackslash}p{2.6cm}>{\raggedright\arraybackslash}X>{\raggedright\arraybackslash}p{2.2cm}cccc>{\raggedright\arraybackslash}p{2.6cm}@{}}
\toprule
Modo de fallo & Efecto & Causa & $S$ & $O$ & $Dt$ & NPR & Medida \\
\midrule
\rule{0pt}{14pt} & & & & & & & \\
\rule{0pt}{14pt} & & & & & & & \\
\rule{0pt}{14pt} & & & & & & & \\
\bottomrule
\end{tabularx}
\end{table}

\section{Registro de decisiones de diseño}

Cada decisión importante (elegir un resolvedor, descartar un modelo, relajar un requisito) se anota con: fecha; contexto; opciones consideradas; decisión; consecuencias y cómo se comprobarán; y quién la tomó. El registro se adjunta al final del documento y la corrección lo usa para valorar la sección~6 y la defensa.

% =====================================================================
\chapter{Autoevaluación}
\label{ap:autoevaluacion}
% =====================================================================

\section{Puntuación por parte y por módulo}

El examen se corrige con la clave del \cref{tab:clave} y con los criterios entre corchetes de cada solución. El \cref{tab:puntos-parte} recoge la puntuación por parte y el \cref{tab:puntos-modulo}, la puntuación por módulo, que sirve para localizar las debilidades: cada módulo vale 20 puntos.

\begin{table}[h]
\centering\small
\caption{Puntuación por parte.}\label{tab:puntos-parte}
\begin{tabular}{@{}lcc@{}}
\toprule
Parte & Puntos & Nota obtenida \\
\midrule
A. Cuestiones (20 de 2 puntos) & 40 & \\
B. Problemas (5 de 12 puntos) & 60 & \\
\midrule
\textbf{Examen} & \textbf{100} & \\
\midrule
C. Proyecto final (se califica aparte) & 100 & \\
\bottomrule
\end{tabular}
\end{table}

\begin{table}[h]
\centering\small
\caption{Puntuación por módulo.}\label{tab:puntos-modulo}
\begin{tabularx}{\textwidth}{@{}>{\raggedright\arraybackslash}X>{\raggedright\arraybackslash}p{2.6cm}>{\raggedright\arraybackslash}p{2.2cm}ccc@{}}
\toprule
Módulo & Cuestiones & Problema & Puntos & Umbral 70\,\% & Obtenidos \\
\midrule
M24. Formulación y resolución & 1--4 (8~p) & B.1 (12~p) & 20 & 14 & \\
M25. Optimización bayesiana y sustitutos & 5--8 (8~p) & B.2 (12~p) & 20 & 14 & \\
M26. Control predictivo y RL & 9--12 (8~p) & B.3 (12~p) & 20 & 14 & \\
M27. Geometría y 3D & 13--16 (8~p) & B.4 (12~p) & 20 & 14 & \\
M28. Búsqueda con LLM, producción y valor & 17--20 (8~p) & B.5 (12~p) & 20 & 14 & \\
\midrule
\textbf{Total} & 20 cuestiones (40~p) & 5 problemas (60~p) & \textbf{100} & \textbf{70} & \\
\bottomrule
\end{tabularx}
\end{table}

El \cref{eje:simp} usa en el apartado~(e) las condiciones de Karush--Kuhn--Tucker de M24, §4.2; se asigna entero a M27 porque el resto del problema es optimización de topología.

\section{Umbral de aprobado}

\begin{itemize}[label=\textendash]
  \item \textbf{Examen de hito}: 70 puntos de 100 (70\,\%) o más.
  \item \textbf{Proyecto final}: 70 puntos de 100 o más en la rúbrica de la opción elegida y todos los criterios de aceptación obligatorios cumplidos (\cref{sec:proy-calificacion}).
  \item \textbf{Puerta de salida de la fase~VI}: además del examen, cada uno de los módulos M24 a M28 exige un test de módulo con al menos 9 respuestas correctas de 12 y una autoevaluación de las prácticas o del laboratorio interactivo de la fase de al menos el 70\,\%.
  \item \textbf{Nota del hito}: $0{,}4\cdot\text{examen}/10+0{,}6\cdot\text{proyecto}/10$, con al menos un 7 en cada parte.
\end{itemize}

Como orientación por módulo, una puntuación inferior a 14 de 20 en el \cref{tab:puntos-modulo} indica que conviene repasar ese módulo aunque el examen esté aprobado. Como la fase~VI cierra el curso, superar este hito completa «Anatomía de un LLM».

\section{Qué hacer si no se alcanza el umbral}

\begin{enumerate}[label=\arabic*.]
  \item Anotar cada fallo en el registro de errores clasificándolo como error de concepto, de cálculo o despiste.
  \item Repasar las secciones del \cref{tab:repaso} correspondientes a los fallos, empezando por los errores de concepto.
  \item Rehacer sin mirar la solución los ejercicios de nivel $\star\star$ y $\star\star\star$ de esos temas en los apuntes del módulo y repetir el caso de estudio correspondiente en su cuaderno de prácticas.
  \item Repetir el test del módulo afectado hasta alcanzar 9 de 12 y volver a presentarse al examen de hito con una versión nueva, no antes de una semana.
  \item Si falla el proyecto, corregir los criterios obligatorios incumplidos con las orientaciones del \cref{ap:correccion} y volver a entregar en un plazo de cuatro semanas; se conserva la nota de los criterios ya superados.
\end{enumerate}

\begin{table}[h]
\centering\small
\caption{Secciones de repaso según el fallo (comprobadas en los índices de los apuntes).}\label{tab:repaso}
\begin{tabularx}{\textwidth}{@{}>{\raggedright\arraybackslash}p{4.2cm}X@{}}
\toprule
Fallo & Repasar \\
\midrule
Cuestiones 1, 2; problema B.1 & M24, §2.1 (programas lineales), §2.2 (dualidad y precios sombra), §2.3 (programación entera mixta, big-M) y §2.4 (ramificación y acotación) \\
Cuestión 3 & M24, §6.2 (vendedor de periódicos) y §6.1 (predecir y después optimizar) \\
Cuestión 4 & M24, §5.4 (optimización multiobjetivo) \\
Cuestiones 5, 8; problema B.2 (c)--(d) & M25, §4.2 (probabilidad de mejora), §4.3 (mejora esperada), §4.4 (cotas de confianza) y §4.5 (varianza latente y predictiva) \\
Cuestión 6; problema B.2 (a), (b), (e) & M25, §3.3 (procesos gaussianos) y §3.4 (hiperparámetros) \\
Cuestión 7 & M25, §6.2 (búsqueda aleatoria frente a rejilla) \\
Cuestión 9 & M26, §1.2 (discretización) y §1.3 (estabilidad) \\
Cuestión 10; problema B.3 (a) & M26, §3.1 (regulador lineal cuadrático) \\
Cuestión 11; problema B.3 (b) & M26, §3.2 (horizonte deslizante) y §3.3 (MPC como programa cuadrático, penalización exacta) \\
Cuestión 12; problema B.3 (c)--(d) & M26, §4.1 (ecuación de Bellman) y §4.4 (Gymnasium: \texttt{terminated} y \texttt{truncated}); M21, §1.3 \\
Cuestiones 13, 14; problema B.4 & M27, §2.1 (flexibilidad mínima), §2.2 (SIMP), §2.3 (método adjunto) y §2.4 (criterios de optimalidad); M24, §4.2 (KKT) \\
Cuestión 15 & M27, §5.3 (error en pantalla) y §5.4 (selección con presupuesto) \\
Cuestión 16 & M27, §1.2 (mallas y característica de Euler) \\
Cuestión 17 & M28, §2.2 (sesgo del ganador) \\
Cuestión 18 & M28, §4.3 (deriva y reentrenamiento) \\
Cuestión 19 & M28, §3.3 (multifidelidad y coste) \\
Cuestión 20; problema B.5 & M28, §5.1 (coste y beneficio), §5.2 (intervalos para la mejora) y §5.3 (caso de estudio); 0.9, §4.4 y §5.4 \\
\bottomrule
\end{tabularx}
\end{table}

\paragraph{Siguiente paso.} Superado el hito, el curso está completo. El proyecto final, con su documento de diseño, su código y su informe, es la pieza que conviene conservar como muestra de trabajo.


\backmatter
% =====================================================================
\begin{thebibliography}{99}
\addcontentsline{toc}{chapter}{Bibliografía}
\markboth{Bibliografía}{}
\small

\subsection*{Exámenes, hojas de problemas y trabajos de curso citados}

\bibitem{mit15053mid}
J.~Orlin y E.~Nasrabadi.
\newblock 15.053: Optimization Methods in Management Science. \emph{Midterm Exam~2: Practice Problem Set} y sus soluciones.
\newblock Massachusetts Institute of Technology, MIT OpenCourseWare, primavera de 2013. Problemas~2 («IP Formulation: Lockboxes») y~4 (restricción «al menos una de dos» y costes no lineales).
\newblock \url{https://ocw.mit.edu/courses/15-053-optimization-methods-in-management-science-spring-2013/pages/study-materials/}

\bibitem{mit15053ps}
J.~Orlin y E.~Nasrabadi.
\newblock 15.053: Optimization Methods in Management Science. Hojas de problemas~4, 5 y~6.
\newblock Massachusetts Institute of Technology, primavera de 2013.
\newblock \url{https://ocw.mit.edu/courses/15-053-optimization-methods-in-management-science-spring-2013/pages/assignments/}

\bibitem{cam4f13}
C.~E. Rasmussen.
\newblock 4F13: Probabilistic Machine Learning. Trabajo de curso~1, «Gaussian processes».
\newblock University of Cambridge, Michaelmas de 2024.
\newblock \url{https://mlg.eng.cam.ac.uk/teaching/4f13/2425/}

\bibitem{ee363final}
S.~Boyd.
\newblock EE363: Linear Dynamical Systems. Examen final (24 horas), invierno de 2006.
\newblock Stanford University. Problema~8, «LQR with linear state cost».
\newblock \url{https://ee363.stanford.edu/archive/notes/final_06.pdf}

\bibitem{ee363}
Stanford University.
\newblock EE363: Linear Dynamical Systems.
\newblock Primavera de 2026. Clases~16 «Linear quadratic regulator: Discrete-time finite horizon», 17 «LQR via Lagrange multipliers» y~18 «Infinite horizon linear quadratic regulator».
\newblock \url{https://ee363.stanford.edu/lectures.html}

\bibitem{mit18s096ps2}
A.~Edelman y S.~G. Johnson.
\newblock 18.S096: Matrix Calculus for Machine Learning and Beyond. Hoja de problemas~2.
\newblock Massachusetts Institute of Technology, MIT OpenCourseWare, IAP (enero) de 2023. Problema~1, gradiente de $f(p)=\bigl(c\T A(p)^{-1}b\bigr)^{2}$ con $A(p)$ tridiagonal mediante una resolución adjunta.
\newblock \url{https://ocw.mit.edu/courses/18-s096-matrix-calculus-for-machine-learning-and-beyond-january-iap-2023/}

\bibitem{mit1805final}
J.~Orloff y J.~Bloom.
\newblock 18.05: Introduction to Probability and Statistics. Examen final.
\newblock Massachusetts Institute of Technology, MIT OpenCourseWare, primavera de 2022. Problemas~9 (contraste $t$ e intervalo de confianza), 11 (potencia) y~13 (intervalo \emph{bootstrap} percentil).
\newblock \url{https://ocw.mit.edu/courses/18-05-introduction-to-probability-and-statistics-spring-2022/}

\bibitem{mit15401}
A.~Lo.
\newblock 15.401: Finance Theory I.
\newblock Massachusetts Institute of Technology, MIT OpenCourseWare, otoño de 2008. Sesiones «Present Value Relations» y «Capital Budgeting».
\newblock \url{https://ocw.mit.edu/courses/15-401-finance-theory-i-fall-2008/}

\subsection*{Libros y artículos}

\bibitem{bt}
D.~Bertsimas y J.~N. Tsitsiklis.
\newblock \emph{Introduction to Linear Optimization}.
\newblock Athena Scientific, 1997. Texto del curso 6.251J del MIT.

\bibitem{bv}
S.~Boyd y L.~Vandenberghe.
\newblock \emph{Convex Optimization}.
\newblock Cambridge University Press, 2004. Texto del curso EE364a de Stanford.
\newblock \url{https://web.stanford.edu/~boyd/cvxbook/}

\bibitem{kw}
M.~J. Kochenderfer y T.~A. Wheeler.
\newblock \emph{Algorithms for Optimization}.
\newblock MIT Press, 2019. Texto del curso AA222/CS361 de Stanford.
\newblock \url{https://algorithmsbook.com/optimization/}

\bibitem{jones}
D.~R. Jones, M.~Schonlau y W.~J. Welch.
\newblock Efficient global optimization of expensive black-box functions.
\newblock \emph{Journal of Global Optimization}, 13:455--492, 1998.

\bibitem{rw}
C.~E. Rasmussen y C.~K.~I. Williams.
\newblock \emph{Gaussian Processes for Machine Learning}.
\newblock MIT Press, 2006.
\newblock \url{https://gaussianprocess.org/gpml/}

\bibitem{garnett}
R.~Garnett.
\newblock \emph{Bayesian Optimization}.
\newblock Cambridge University Press, 2023.
\newblock \url{https://bayesoptbook.com/}

\bibitem{rmd}
J.~B. Rawlings, D.~Q. Mayne y M.~M. Diehl.
\newblock \emph{Model Predictive Control: Theory, Computation, and Design}, 2.ª ed.
\newblock Nob Hill Publishing, 2017. Caps.~1 y~2.
\newblock \url{https://sites.engineering.ucsb.edu/~jbraw/mpc/}

\bibitem{sb}
R.~S. Sutton y A.~G. Barto.
\newblock \emph{Reinforcement Learning: An Introduction}, 2.ª ed.
\newblock MIT Press, 2018.
\newblock \url{http://incompleteideas.net/book/the-book-2nd.html}

\bibitem{bs}
M.~P. Bendsøe y O.~Sigmund.
\newblock \emph{Topology Optimization: Theory, Methods, and Applications}, 2.ª ed.
\newblock Springer, Berlín, 2004. Cap.~1 «Topology optimization by distribution of isotropic material».

\bibitem{andreassen}
E.~Andreassen, A.~Clausen, M.~Schevenels, B.~S. Lazarov y O.~Sigmund.
\newblock Efficient topology optimization in MATLAB using 88 lines of code.
\newblock \emph{Structural and Multidisciplinary Optimization}, 43(1):1--16, 2011.

\bibitem{garland}
M.~Garland y P.~S. Heckbert.
\newblock Surface simplification using quadric error metrics.
\newblock En \emph{Proceedings of SIGGRAPH 97}, ACM, 1997.

\bibitem{bp}
R.~E. Barlow y F.~Proschan.
\newblock \emph{Mathematical Theory of Reliability}.
\newblock SIAM, 1996 (original: Wiley, 1965). Cap.~4, «Optimum Maintenance Policies».

\bibitem{montgomery}
D.~C. Montgomery.
\newblock \emph{Design and Analysis of Experiments}, 9.ª ed.
\newblock Wiley, 2017. Cap.~2 «Simple Comparative Experiments».

\subsection*{Cursos de ingeniería de sistemas y de sistemas con aprendizaje automático}

\bibitem{mit16842}
O.~de Weck.
\newblock 16.842: Fundamentals of Systems Engineering.
\newblock Massachusetts Institute of Technology, MIT OpenCourseWare, otoño de 2015.
\newblock \url{https://ocw.mit.edu/courses/16-842-fundamentals-of-systems-engineering-fall-2015/}

\bibitem{cs329s}
C.~Huyen y equipo docente de CS329S.
\newblock CS329S: Machine Learning Systems Design.
\newblock Stanford University, invierno de 2022.
\newblock \url{https://stanford-cs329s.github.io/syllabus.html}

\bibitem{huyen}
C.~Huyen.
\newblock \emph{Designing Machine Learning Systems}.
\newblock O'Reilly, 2022.

\bibitem{sculley}
D.~Sculley, G.~Holt, D.~Golovin, E.~Davydov, T.~Phillips, D.~Ebner, V.~Chaudhary, M.~Young, J.-F. Crespo y D.~Dennison.
\newblock Hidden technical debt in machine learning systems.
\newblock En \emph{Advances in Neural Information Processing Systems 28 (NeurIPS)}, 2015.

\subsection*{Apuntes del curso}

\bibitem{apuntes}
\emph{Anatomía de un LLM}. Apuntes de los módulos~24 a~28 (fase~VI) y de los módulos de las fases anteriores citados.
\newblock Edición de octubre de 2026.

\end{thebibliography}


\end{document}
